% Blueprint content for the One ID rules.
%
% Do not write \leanok here.  scripts/render_status.py reads build/status.json
% and writes content.generated.tex with \leanok added to every environment
% whose declaration the build reports as defined or proven.
%
% A node with \lean{} is a Lean declaration.  A node without it is a candidate
% from the survey in docs/one-id-survey.md, with its ID in the title; one the
% guide states vaguely or not at all is marked \notready.  The progress page
% reads the paragraphs that start with Counterexample:, Open point:, Blocks:,
% External: and Source:.

\chapter{One ID credentials and their presentation}

This blueprint follows IATA's One ID Technical Implementation Guide, version
26.3: the credentials a passenger's wallet holds, the photoid and the boarding
pass, and how a wallet presents them to a relying party.  A node with a Lean
name is a definition or a theorem in the model.  A green node is defined or
proven in Lean; an uncolored theorem is open, which means the guide as written
does not give the guarantee, and an example the build computes shows the case.
The other nodes are candidate rules from a survey of every rule in the guide.
They have no Lean yet.  A candidate the guide states vaguely or not at all is
marked not ready, and its open point says what has to be decided.

The guide is cited by section number.  It builds on ISO, OpenID, W3C, IETF,
ICAO and IATA specifications that this project has not read; a node that
depends on one names it under External.

The review document, under \texttt{review/} on the site, shows the Lean
declarations with their docstrings, the examples the build computes, and the
open questions.

\section{Core concepts}

\begin{definition}[Instant]
  \label{def:instant}
  \lean{Rules.Core.Instant}
  A point in time, as whole seconds since the Unix epoch.  All times are UTC.
\end{definition}

\begin{definition}[Credential]
  \label{def:mdoc}
  \lean{Rules.OneId.Mdoc, Rules.OneId.NameSpace, Rules.OneId.DataElement, Rules.OneId.Value}
  \uses{def:instant}
  A digital credential in the ISO mdoc format, as far as the guide describes
  one: a document type, namespaces of data elements, and the instant its
  validity ends.  The guide uses claim, data element and data attribute for
  the same thing.
\end{definition}

\begin{definition}[Findings and recommendations]
  \label{def:finding}
  \lean{Rules.OneId.Finding, Rules.OneId.Advisory}
  \uses{def:mdoc}
  A finding is a reason a credential breaks something the guide says with
  MUST or SHALL.  An advisory is something the guide says with SHOULD or
  RECOMMENDED that a credential does not follow; it does not make the
  credential unacceptable.
\end{definition}

% The credential data model (§4.1): rules and guarantees that are in Lean.
% Written by hand.  Do not write \leanok; scripts/render_status.py adds it.
% The candidates of this area that have no Lean yet are in
% credentials-candidates.tex.

\section{Credential data model}

\begin{definition}[Data elements and their names (C-25)]
  \label{def:element}
  \lean{Rules.OneId.Mdoc.element?, Rules.OneId.Mdoc.wellFormed, Rules.OneId.Mdoc.nameSpace?, Rules.OneId.NameSpace.element?}
  \uses{def:mdoc}
  Each data element of a credential is identified by the name of its
  namespace and its identifier.  The names of a credential are unique when no
  two namespaces share a name and no two data elements of one namespace share
  an identifier.  The guide leaves the structure to ISO/IEC 18013-5, which
  this project has not read, and does not say that names are unique.

  External: ISO/IEC 18013-5:2021 §6.3.1, §8.3.1, §8.3.2.

  Source: \href{https://standards.atlassian.net/wiki/spaces/OneID/pages/2947416082/Technical+Implementation+Guide}{One ID guide} §4.1, §4.1.1.
\end{definition}

\begin{theorem}[Every data element can be requested on its own]
  \label{thm:element-addressable}
  \lean{Rules.OneId.element_addressable}
  \uses{def:element}
  In a credential whose names are unique, every data element is found by the
  name of its namespace and its identifier.
\end{theorem}
\begin{proof}
  \uses{def:element}
  With unique names, the search for a name finds the element that carries it.
\end{proof}

\begin{theorem}[No value for a name the credential lacks]
  \label{thm:element-absent}
  \lean{Rules.OneId.element_absent}
  \uses{def:element}
  A credential returns no value for a namespace name and an identifier when
  none of its namespaces of that name holds a data element with that
  identifier.
\end{theorem}
\begin{proof}
  \uses{def:element}
  The search for the namespace or for the element finds nothing.
\end{proof}

\begin{theorem}[The names in a credential are unique]
  \label{thm:names-unique}
  \lean{Rules.OneId.photoid_names_unique, Rules.OneId.boardingPass_names_unique}
  \uses{def:element, def:photoid, def:boarding-pass}
  The names in an accepted photoid, and in an accepted boarding pass, are
  unique.  Not proven.

  Counterexample: a photoid whose base identity namespace holds two family
  names, SMITH and JONES, is accepted under the reading that names may
  repeat.  A relying party that requests the family name receives one of the
  two.  A boarding pass with two boarding pass strings is accepted in the same
  way.

  Open point: the guide does not say that namespace names and data element
  identifiers are unique.  May a name repeat within a credential, and is such
  a credential rejected?

  Lookup: ISO/IEC 18013-5:2021 §8.3.1, §8.3.2.

  Source: \href{https://standards.atlassian.net/wiki/spaces/OneID/pages/2947416082/Technical+Implementation+Guide}{One ID guide} §4.1.
\end{theorem}

\begin{definition}[Photoid requirements (C-27)]
  \label{def:photoid}
  \lean{Rules.OneId.Photoid.violations, Rules.OneId.PhotoidReading, Rules.OneId.PhotoidTables, Rules.OneId.ElementTable, Rules.OneId.TableEntry, Rules.OneId.Presence, Rules.OneId.Encoding}
  \uses{def:element, def:finding}
  A credential is accepted as a photoid only when it has the document type
  \texttt{org.iso.23220.photoid.1} and the base identity namespace
  \texttt{org.iso.23220.1}, with every data element Table C.1 of ISO/IEC TS
  23220-4 marks mandatory and each listed element in the encoding the table
  gives.  The tables are a parameter, since this project has not read them.
  The points the guide leaves undecided are a reading: whether the portrait
  is mandatory, whether data the guide does not list is allowed, whether the
  portrait must be the same as data group 2, whether a recommended namespace
  must be complete once it is present, and whether names may repeat.

  External: ISO/IEC TS 23220-4:2026 Annex C, §C.2.1 and Tables C.1 to C.3;
  ISO/IEC TS 23220-2:2026.

  Source: \href{https://standards.atlassian.net/wiki/spaces/OneID/pages/2947416082/Technical+Implementation+Guide}{One ID guide} §4.1.1, §2.
\end{definition}

\begin{theorem}[A photoid has the photoid document type]
  \label{thm:photoid-doctype}
  \lean{Rules.OneId.photoid_docType}
  \uses{def:photoid}
  Under every reading and every content of the tables, an accepted photoid
  has the photoid document type.
\end{theorem}
\begin{proof}
  \uses{def:photoid}
  Another document type is a finding.
\end{proof}

\begin{theorem}[A photoid has the base identity namespace]
  \label{thm:photoid-has-base}
  \lean{Rules.OneId.photoid_has_base}
  \uses{def:photoid}
  Under every reading and every content of the tables, an accepted photoid
  has the base identity namespace.
\end{theorem}
\begin{proof}
  \uses{def:photoid}
  A missing namespace is a finding.
\end{proof}

\begin{theorem}[A photoid has every mandatory base element]
  \label{thm:photoid-mandatory-base}
  \lean{Rules.OneId.photoid_has_mandatory_base}
  \uses{def:photoid}
  An accepted photoid holds every data element Table C.1 marks mandatory in
  the base identity namespace, whatever the table lists.
\end{theorem}
\begin{proof}
  \uses{def:photoid}
  Each mandatory element that is missing is a finding.
\end{proof}

\begin{theorem}[Base elements have the encoding of their table]
  \label{thm:photoid-base-encoding}
  \lean{Rules.OneId.photoid_base_encoding}
  \uses{def:photoid}
  In an accepted photoid, every data element of the base identity namespace
  that Table C.1 lists has the encoding the table gives it, whatever the
  table says.  The model tells a text string, a byte string and a yes or no
  value apart.
\end{theorem}
\begin{proof}
  \uses{def:photoid}
  A listed element in another encoding is a finding.
\end{proof}

\begin{theorem}[Every photoid has a portrait]
  \label{thm:photoid-has-portrait}
  \lean{Rules.OneId.photoid_has_portrait}
  \uses{def:photoid, def:c04, def:c103}
  Every accepted photoid has a portrait.  Not proven.

  Counterexample: a photoid with a family name and a given name and no
  portrait is accepted when Table C.1 does not mark the portrait mandatory and the
  guide's sentence is read as a rule about the portrait's content only.

  Open point: the guide says the portrait MUST hold the passport's image and
  that the business requirements need it, and it also refers the presence of
  every base element to Table C.1.  Is the portrait mandatory in a One ID
  photoid, whatever Table C.1 says?

  Lookup: ISO/IEC TS 23220-4:2026 Annex C, Table C.1.  If the table marks the
  portrait mandatory, the guarantee follows.

  Source: \href{https://standards.atlassian.net/wiki/spaces/OneID/pages/2947416082/Technical+Implementation+Guide}{One ID guide} §4.1.1.
\end{theorem}

\begin{theorem}[A photoid carries only what the guide lists]
  \label{thm:photoid-only-listed}
  \lean{Rules.OneId.photoid_only_listed}
  \uses{def:photoid}
  An accepted photoid carries only the three namespaces the guide lists, and
  in each only data elements its table lists.  Not proven.

  Counterexample: a photoid that carries a fourth namespace with a frequent
  flyer number, and a data element in its base identity namespace that Table
  C.1 does not list, is accepted under the reading that the lists are not
  closed.

  Open point: the guide lists three namespaces and refers to a table for
  each, and does not say the lists are closed.  May a One ID photoid carry a
  namespace or a data element the guide and its tables do not list?

  Source: \href{https://standards.atlassian.net/wiki/spaces/OneID/pages/2947416082/Technical+Implementation+Guide}{One ID guide} §4.1.1.
\end{theorem}

\begin{theorem}[The portrait is the same as data group 2]
  \label{thm:portrait-is-dg2}
  \lean{Rules.OneId.portrait_is_dg2}
  \uses{def:photoid}
  When an accepted photoid carries both a portrait and a data group 2, they
  are the same.  Not proven.

  Counterexample: a photoid whose portrait is one face image and whose data
  group 2 is another is accepted under the reading that ``equivalent'' is
  left to the issuer.

  Open point: the guide says the portrait MUST hold ``the dg2 or equivalent
  image data'' and does not say what equivalent means or who checks it.  When
  a photoid carries data group 2, must the portrait be that image, and must a
  relying party check it?

  Source: \href{https://standards.atlassian.net/wiki/spaces/OneID/pages/2947416082/Technical+Implementation+Guide}{One ID guide} §4.1.1.
\end{theorem}

\begin{theorem}[A photoid namespace that is present is complete]
  \label{thm:photoid-namespace-complete}
  \lean{Rules.OneId.photoid_nameSpace_complete}
  \uses{def:photoid, def:c103}
  When an accepted photoid carries the photoid namespace, the namespace holds
  every data element Table C.2 marks mandatory.  Not proven.

  Counterexample: a photoid whose photoid namespace holds only the image
  captured at enrolment is accepted under the reading that the table is only
  recommended, though the table marks a document number mandatory.

  Open point: the guide recommends ``this namespace and data elements
  presence'' in one sentence.  Once the photoid namespace is present, are the
  data elements Table C.2 marks mandatory required?

  Source: \href{https://standards.atlassian.net/wiki/spaces/OneID/pages/2947416082/Technical+Implementation+Guide}{One ID guide} §4.1.1.
\end{theorem}

\begin{theorem}[A data groups namespace that is present is complete]
  \label{thm:datagroups-namespace-complete}
  \lean{Rules.OneId.dataGroups_nameSpace_complete}
  \uses{def:photoid, def:c103}
  When an accepted photoid carries the data groups namespace, the namespace
  holds every data element Table C.3 marks mandatory.  Not proven.

  Counterexample: a photoid whose data groups namespace holds only data group
  2 is accepted under the reading that the table is only recommended, though
  the table marks the passport's security object mandatory.

  Open point: the guide recommends the data groups namespace and the presence
  of its data elements as in Table C.3 in one sentence.  Once the data groups
  namespace is present, are the data elements Table C.3 marks mandatory
  required?

  Source: \href{https://standards.atlassian.net/wiki/spaces/OneID/pages/2947416082/Technical+Implementation+Guide}{One ID guide} §4.1.1.
\end{theorem}

\begin{definition}[Source document of a photoid (C-06)]
  \label{def:photoid-source}
  \lean{Rules.OneId.Photoid.sourceFindings, Rules.OneId.SourceDocument}
  \uses{def:photoid}
  The kind of identity document a photoid was derived from is a fact the
  issuer knows.  Whether a source other than a passport is a reason to reject
  the photoid is a reading: the terminology allows an equivalent identity
  document, and §4.1.1 says the portrait MUST hold image data sourced from
  the passport.

  Source: \href{https://standards.atlassian.net/wiki/spaces/OneID/pages/2947416082/Technical+Implementation+Guide}{One ID guide} §2, §4.1.1.
\end{definition}

\begin{theorem}[A photoid comes from a passport]
  \label{thm:photoid-from-passport}
  \lean{Rules.OneId.photoid_from_passport}
  \uses{def:photoid-source}
  A photoid the system accepts was derived from a passport.  Not proven.

  Counterexample: a photoid derived from a national identity card is accepted
  under the reading that an equivalent identity document is allowed.  With
  the photoid namespace absent, nothing in the credential says which document
  it came from.

  Open point: the terminology says a photoid can be derived from a passport
  or an equivalent identity document, and §4.1.1 says its values should
  always come from the passport and its portrait MUST hold image data sourced
  from the passport.  May a One ID photoid be derived from a document other
  than a passport?

  Source: \href{https://standards.atlassian.net/wiki/spaces/OneID/pages/2947416082/Technical+Implementation+Guide}{One ID guide} §2, §4.1.1.
\end{theorem}

\begin{definition}[Photoid recommendations (C-32)]
  \label{def:photoid-advisories}
  \lean{Rules.OneId.Photoid.advisories}
  \uses{def:photoid}
  A photoid that meets every requirement is accepted, and is told which
  recommendations of the guide it does not follow: the photoid namespace with
  the travel document's data, the image captured at enrolment, the data
  groups namespace with the passport chip's data, and values derived from a
  passport.

  Open point: without the photoid namespace and the data groups namespace,
  nothing in a photoid says which document it was derived from, and a relying
  party cannot check the portrait against the passport chip.  Is that
  acceptable for the admissibility use case?

  External: ISO/IEC TS 23220-4:2026 Annex C, Tables C.2 and C.3; ICAO Doc 9303
  Part 10.

  Source: \href{https://standards.atlassian.net/wiki/spaces/OneID/pages/2947416082/Technical+Implementation+Guide}{One ID guide} §4.1.1, §2.
\end{definition}

\begin{theorem}[A missing photoid namespace is reported]
  \label{thm:advises-photoid-namespace}
  \lean{Rules.OneId.advises_photoid_nameSpace}
  \uses{def:photoid-advisories}
  A photoid that lacks the photoid namespace is told so, as a recommendation
  it does not follow.
\end{theorem}
\begin{proof}
  \uses{def:photoid-advisories}
  By the definition.
\end{proof}

\begin{theorem}[A missing data groups namespace is reported]
  \label{thm:advises-datagroups-namespace}
  \lean{Rules.OneId.advises_dataGroups_nameSpace}
  \uses{def:photoid-advisories}
  A photoid that lacks the data groups namespace is told so, as a
  recommendation it does not follow.
\end{theorem}
\begin{proof}
  \uses{def:photoid-advisories}
  By the definition.
\end{proof}

\begin{theorem}[A missing enrolment image is reported]
  \label{thm:advises-enrolment-portrait}
  \lean{Rules.OneId.advises_enrolment_portrait}
  \uses{def:photoid-advisories}
  A photoid whose photoid namespace lacks the image captured at enrolment is
  told so, whatever Table C.2 says about that image.
\end{theorem}
\begin{proof}
  \uses{def:photoid-advisories}
  By cases on whether the table marks the image mandatory.
\end{proof}

\begin{theorem}[A source other than a passport is reported]
  \label{thm:advises-source-document}
  \lean{Rules.OneId.advises_source_document}
  \uses{def:photoid-advisories}
  A photoid derived from a document other than a passport is told so.
\end{theorem}
\begin{proof}
  \uses{def:photoid-advisories}
  By the definition.
\end{proof}

\begin{theorem}[A missing recommended namespace is not a reason to reject]
  \label{thm:absent-recommended-namespace}
  \lean{Rules.OneId.recommended_nameSpace_never_required}
  \uses{def:photoid-advisories, def:photoid}
  No photoid is rejected for lacking the photoid namespace or the data groups
  namespace: neither is among the requirements a photoid can break, under
  every reading and every content of the tables.
\end{theorem}
\begin{proof}
  \uses{def:photoid}
  No part of the check reports either namespace as missing.
\end{proof}

\begin{definition}[Boarding pass requirements (C-36)]
  \label{def:boarding-pass}
  \lean{Rules.OneId.BoardingPass.violations, Rules.OneId.BoardingPassReading, Rules.OneId.Resolution792, Rules.OneId.BoardingPass.bcbpString}
  \uses{def:element, def:finding}
  A credential is accepted as a boarding pass only when it has the document
  type \texttt{org.iata.boardingpass.1}, unique names, and the namespace
  \texttt{org.iata.bcbp.1} with the data element \texttt{bcbp\_string}, a text
  string that follows IATA Resolution 792.  What the resolution decides is a
  parameter, since this project has not read it.  Whether a boarding pass may
  carry anything else, and whether names may repeat, is a reading.

  External: IATA Resolution 792; IETF RFC 8949.

  Source: \href{https://standards.atlassian.net/wiki/spaces/OneID/pages/2947416082/Technical+Implementation+Guide}{One ID guide} §4.1.2.
\end{definition}

\begin{theorem}[A boarding pass has the boarding pass document type]
  \label{thm:bp-doctype}
  \lean{Rules.OneId.boardingPass_docType}
  \uses{def:boarding-pass}
  Under every reading, an accepted boarding pass has the boarding pass
  document type.
\end{theorem}
\begin{proof}
  \uses{def:boarding-pass}
  Another document type is a finding.
\end{proof}

\begin{theorem}[A boarding pass holds one boarding pass string]
  \label{thm:bp-has-string}
  \lean{Rules.OneId.boardingPass_has_string}
  \uses{def:boarding-pass}
  An accepted boarding pass holds a boarding pass string: its namespace and
  the data element are present, the value is a text string, and the string
  follows Resolution 792, whatever the resolution says.
\end{theorem}
\begin{proof}
  \uses{def:boarding-pass}
  Each of the four ways the string can be wrong is a finding.
\end{proof}

\begin{theorem}[No credential is both a photoid and a boarding pass]
  \label{thm:not-both}
  \lean{Rules.OneId.not_photoid_and_boardingPass}
  \uses{def:photoid, def:boarding-pass, thm:photoid-doctype, thm:bp-doctype}
  No credential is accepted both as a photoid and as a boarding pass.
\end{theorem}
\begin{proof}
  \uses{thm:photoid-doctype, thm:bp-doctype}
  The two document types differ.
\end{proof}

\begin{theorem}[A boarding pass carries only what the guide lists]
  \label{thm:bp-only-listed}
  \lean{Rules.OneId.boardingPass_only_listed}
  \uses{def:boarding-pass}
  An accepted boarding pass carries one namespace with one data element, the
  boarding pass string.  Not proven.

  Counterexample: a boarding pass that carries a second namespace with the
  seat as a data element of its own is accepted under the reading that more
  than the guide lists is allowed.

  Open point: the guide lists one namespace with one data element and says a
  later version may add a namespace.  May a boarding pass credential carry a
  namespace or a data element the guide does not list today?

  Source: \href{https://standards.atlassian.net/wiki/spaces/OneID/pages/2947416082/Technical+Implementation+Guide}{One ID guide} §4.1.2.
\end{theorem}

\begin{definition}[Expiry of a credential (C-44)]
  \label{def:expired}
  \lean{Rules.OneId.expired, Rules.OneId.ExpiryReading}
  \uses{def:mdoc, def:instant}
  A credential is expired when the current time is after the end of the
  validity its issuer signed.  Whether it is still valid at the exact instant
  is a reading: the check is in ISO/IEC 18013-5, which this project has not
  read.

  External: ISO/IEC 18013-5:2021 §9.3.1.

  Source: \href{https://standards.atlassian.net/wiki/spaces/OneID/pages/2947416082/Technical+Implementation+Guide}{One ID guide} §4.1.2, §4.2.1.2.
\end{definition}

\begin{theorem}[Expired after the end of validity]
  \label{thm:expired-after}
  \lean{Rules.OneId.expired_after_validUntil}
  \uses{def:expired}
  A credential is expired once the current time is after the end of its
  validity, under both readings of the exact instant.
\end{theorem}
\begin{proof}
  \uses{def:expired}
  By the definition, in each reading.
\end{proof}

\begin{theorem}[Valid before the end of validity]
  \label{thm:valid-before}
  \lean{Rules.OneId.valid_before_validUntil}
  \uses{def:expired}
  A credential is not expired while the current time is before the end of its
  validity, under both readings of the exact instant.
\end{theorem}
\begin{proof}
  \uses{def:expired}
  By the definition, in each reading.
\end{proof}

\begin{theorem}[An expired credential stays expired]
  \label{thm:expired-stays}
  \lean{Rules.OneId.expired_stays_expired}
  \uses{def:expired}
  Once a credential has expired it is expired at every later time.
\end{theorem}
\begin{proof}
  \uses{def:expired}
  Time only moves further past the end of validity.
\end{proof}

\begin{theorem}[Valid at the end of validity]
  \label{thm:valid-at-validuntil}
  \lean{Rules.OneId.valid_at_validUntil}
  \uses{def:expired}
  A credential is not expired at the exact instant its validity ends.  Not
  proven.

  Counterexample: a boarding pass whose validity ends exactly at departure,
  which the guide allows, is expired at the moment of departure under the
  reading that validity ends at the instant itself.

  Open point: is a credential still valid at the exact instant its validity
  ends?

  Lookup: ISO/IEC 18013-5:2021 §9.3.1.

  Source: \href{https://standards.atlassian.net/wiki/spaces/OneID/pages/2947416082/Technical+Implementation+Guide}{One ID guide} §4.1.2, §4.2.1.2.
\end{theorem}

\begin{definition}[Validity of a boarding pass (C-38)]
  \label{def:bp-validity}
  \lean{Rules.OneId.BoardingPass.validityFindings, Rules.OneId.ValidityReading, Rules.OneId.DepartureReading, Rules.OneId.Journey, Rules.OneId.Journey.departure}
  \uses{def:boarding-pass, def:expired}
  The validity of a boarding pass is accepted only if it ends on or soon
  after the departure of the last leg of the journey.  Which departure is
  meant, the scheduled one or the actual one, and how long ``soon after''
  lasts are a reading: the guide says neither.

  External: ISO/IEC 18013-5:2021, validity information.

  Source: \href{https://standards.atlassian.net/wiki/spaces/OneID/pages/2947416082/Technical+Implementation+Guide}{One ID guide} §4.1.2, implementation consideration 1.
\end{definition}

\begin{theorem}[Validity reaches the departure]
  \label{thm:validity-reaches-departure}
  \lean{Rules.OneId.validity_reaches_departure}
  \uses{def:bp-validity}
  The validity of an accepted boarding pass does not end before the departure
  of the last leg, whichever departure the reading names and however long
  ``soon after'' is.
\end{theorem}
\begin{proof}
  \uses{def:bp-validity}
  An earlier end is a finding.
\end{proof}

\begin{theorem}[Validity ends soon after the departure]
  \label{thm:validity-ends-soon-after}
  \lean{Rules.OneId.validity_ends_soon_after}
  \uses{def:bp-validity}
  The validity of an accepted boarding pass ends no later than ``soon after''
  the departure of the last leg, for the departure and the length the reading
  gives.
\end{theorem}
\begin{proof}
  \uses{def:bp-validity}
  A later end is a finding.
\end{proof}

\begin{theorem}[Valid before the departure]
  \label{thm:valid-before-departure}
  \lean{Rules.OneId.valid_before_departure}
  \uses{def:bp-validity, thm:valid-before, thm:validity-reaches-departure}
  An accepted boarding pass is not expired before the departure of the last
  leg, for the departure the reading names.
\end{theorem}
\begin{proof}
  \uses{thm:valid-before, thm:validity-reaches-departure}
  Its validity reaches the departure, and a credential is valid before the
  end of its validity.
\end{proof}

\begin{theorem}[Expired once the journey is completed]
  \label{thm:expired-once-journey-completed}
  \lean{Rules.OneId.expired_once_journey_completed}
  \uses{def:bp-validity, def:c42}
  An accepted boarding pass is expired once the journey is completed.  Not
  proven.

  Counterexample: a flight leaves at 10:00 and lands at 12:00.  With ``soon
  after'' read as a day, a boarding pass whose validity ends at 10:00 the next
  day is accepted, and at 14:00 it is still valid.

  Open point: the guide says the end of validity makes the boarding pass
  expire once the journey is completed, gives no limit for ``soon after'' and
  no definition of a completed journey, and puts revocation out of scope.  How
  long after the last departure may the validity of a boarding pass run?

  Source: \href{https://standards.atlassian.net/wiki/spaces/OneID/pages/2947416082/Technical+Implementation+Guide}{One ID guide} §4.1.2, implementation consideration 1.
\end{theorem}

\begin{theorem}[Valid until the flight actually departs]
  \label{thm:valid-until-actual-departure}
  \lean{Rules.OneId.valid_until_actual_departure}
  \uses{def:bp-validity}
  An accepted boarding pass is not expired before its last leg actually
  departs.  Not proven.

  Counterexample: a flight scheduled for 10:00 departs at 13:00.  With the
  scheduled departure meant and ``soon after'' read as an hour, a boarding
  pass whose validity ends at 11:00 is accepted, and at 12:30 it is expired:
  the passenger cannot present it to board.

  Open point: an issuer knows only the scheduled departure when it signs.
  Which departure is meant, and what happens to a boarding pass credential
  when its flight is delayed past the end of its validity?

  Source: \href{https://standards.atlassian.net/wiki/spaces/OneID/pages/2947416082/Technical+Implementation+Guide}{One ID guide} §4.1.2, implementation consideration 1.
\end{theorem}

\begin{definition}[Validity of a photoid (C-35)]
  \label{def:photoid-validity}
  \lean{Rules.OneId.Photoid.validityFindings, Rules.OneId.PhotoidValidityReading}
  \uses{def:photoid, def:expired}
  The guide gives no rule for the validity of a photoid.  Two readings are
  kept: the issuer may set any end of validity, or the validity ends no later
  than the document the photoid was derived from expires.

  Source: \href{https://standards.atlassian.net/wiki/spaces/OneID/pages/2947416082/Technical+Implementation+Guide}{One ID guide} §4.1.1, §1.4.
\end{definition}

\begin{theorem}[A photoid expires with its passport]
  \label{thm:photoid-expires-with-document}
  \lean{Rules.OneId.photoid_expires_with_document}
  \uses{def:photoid-validity}
  A photoid whose validity was accepted is expired once the document it was
  derived from has expired.  Not proven.

  Counterexample: a photoid whose validity ends on 1 January 2030, derived
  from a passport that expires on 30 June 2027, is accepted under the reading
  that the issuer may set any end of validity.  Two weeks after the passport
  expired the photoid is still valid.

  Open point: must the validity of a photoid end no later than the expiry of
  the passport it was derived from?

  Source: \href{https://standards.atlassian.net/wiki/spaces/OneID/pages/2947416082/Technical+Implementation+Guide}{One ID guide} §4.1.1, §1.4.
\end{theorem}

\begin{definition}[One traveler for a photoid and a boarding pass (C-41)]
  \label{def:traveler-match}
  \lean{Rules.OneId.travelerMatch, Rules.OneId.MatchReading, Rules.OneId.Issued, Rules.OneId.Traveler, Rules.OneId.Photoid.name}
  \uses{def:photoid, def:boarding-pass}
  A photoid and a boarding pass presented together are accepted as a pair
  only if both are about the same traveler.  The guide states no rule for
  deciding it, so the check is a reading: nothing is checked, which is the
  guide as written, or the names in the two credentials agree.  Other
  readings, such as one that uses the device key, are added as the rules they
  need are modelled.

  External: IATA Resolution 792; ISO/IEC TS 23220-4:2026 Annex C, Table C.1.

  Source: \href{https://standards.atlassian.net/wiki/spaces/OneID/pages/2947416082/Technical+Implementation+Guide}{One ID guide} §4.1.2, §3.2.
\end{definition}

\begin{theorem}[An accepted pair is one traveler's]
  \label{thm:matched-pair-same-traveler}
  \lean{Rules.OneId.matched_pair_same_traveler}
  \uses{def:traveler-match, def:c03, def:c19, def:c42}
  A photoid and a boarding pass that are accepted as a pair are about the
  same traveler.  Not proven.

  Counterexample: with no check between the two, Alice's photoid is accepted
  with Bruno's boarding pass.  With the names compared, the photoid of one
  John Smith is accepted with the boarding pass of another John Smith.

  Open point: the guide requires a boarding pass together with the face image
  of a photoid at a contactless touchpoint and names no data element to
  compare.  No check of the two credentials' content settles it by itself,
  since who a credential is about is fixed at issuance, which the guide puts
  out of scope.  What ties a photoid and a boarding pass to one traveler: a
  rule for issuers, the device key, or something else, and what must a
  relying party check?

  Source: \href{https://standards.atlassian.net/wiki/spaces/OneID/pages/2947416082/Technical+Implementation+Guide}{One ID guide} §4.1.2, §3.2.
\end{theorem}

\begin{theorem}[No check of content alone identifies the traveler]
  \label{thm:no-content-check}
  \lean{Rules.OneId.no_content_check_identifies_the_traveler}
  \uses{def:traveler-match}
  Whatever a relying party checks on the content of a photoid and a boarding
  pass, a pair it accepts can be about two different travelers: the same two
  credentials could have been issued about different people.  So a matching
  rule gives one traveler only together with a rule about issuance.
\end{theorem}
\begin{proof}
  \uses{def:traveler-match}
  The subject of a credential is not part of its content.
\end{proof}

\begin{theorem}[A traveler's own pair is accepted]
  \label{thm:same-traveler-matches}
  \lean{Rules.OneId.same_traveler_matches}
  \uses{def:traveler-match, def:c102}
  A photoid and a boarding pass about the same traveler are accepted as a
  pair.  Not proven.

  Counterexample: a traveler whose passport gives GARCÍA and JOSÉ MARÍA, and
  whose boarding pass carries GARCIA and JOSEMARIA, is refused under an exact
  comparison of names.

  Open point: if names are compared, how are a passport name and a boarding
  pass name brought to the same form?

  Source: \href{https://standards.atlassian.net/wiki/spaces/OneID/pages/2947416082/Technical+Implementation+Guide}{One ID guide} §4.1.2.
\end{theorem}

% Candidate rules of the credential data model that have no Lean yet (§4.1).
% Written by hand; see docs/one-id-survey.md.  The section heading is in credentials.tex.

\begin{definition}[The security model of the photoid (C-26)]
  \label{def:c26}
  \uses{def:element, def:photoid}
  The security model of ISO/IEC TS 23220-4 Annexes A and B must be followed:
  an issuer-signed security object, device-signed responses, and selective
  disclosure of individual data elements. The guide says these ensure a
  verifier receives only what it requested and the holder consented to, with
  authenticity and integrity.

  The sentence states guarantee G-01 as a consequence of the ISO mechanisms.
  The model has to derive it from the checks the guide lists, not assume it.

  External: ISO/IEC TS 23220-4:2026 Annex A, Annex B.

  Source:
  \href{https://standards.atlassian.net/wiki/spaces/OneID/pages/2947416082/Technical+Implementation+Guide}{One
  ID guide} §4.1.1.
\end{definition}

\begin{definition}[Age attestation (C-31)]
  \label{def:c31}
  \uses{def:element, def:photoid, def:c04, def:c103}
  The base identity namespace includes age attestation elements, which
  disclose an age limit without the age or the date of birth when a relying
  party requests them alone.

  An age attestation is a data element like any other, so the rule for
  addressing one element covers it. The guarantee that the date of birth is
  not disclosed with it belongs to G-01.

  External: ISO/IEC TS 23220-4:2026 Annex C, Table C.1.

  Source:
  \href{https://standards.atlassian.net/wiki/spaces/OneID/pages/2947416082/Technical+Implementation+Guide}{One
  ID guide} §4.1.1.
\end{definition}

\begin{definition}[Showing the boarding pass in the wallet (C-39)]
  \label{def:c39}
  \uses{def:boarding-pass, def:c05, def:c102}
  \notready
  The wallet must display 'only all mandatory elements' of the boarding pass
  string in a user-friendly way.

  Needs the list of mandatory elements from Resolution 792. Not modelled.

  Open point: Does the wallet show all mandatory elements and nothing else,
  or at least the mandatory ones?

  External: IATA Resolution 792.

  Source:
  \href{https://standards.atlassian.net/wiki/spaces/OneID/pages/2947416082/Technical+Implementation+Guide}{One
  ID guide} §4.1.2, consideration 2.
\end{definition}

\begin{definition}[How many boarding pass credentials a journey has (C-42)]
  \label{def:c42}
  \uses{def:boarding-pass, def:c03, def:c102}
  \notready
  The guide speaks of the last leg of a journey, so one boarding pass
  credential may cover several flights. It does not say whether a journey has
  one credential or one per leg, or what a wallet holds for several
  passengers.

  A request names a document type. If a wallet holds several boarding passes,
  the rule that picks the one to send is missing.

  Open point: Is there one boarding pass credential per journey, per flight
  or per boarding pass string? When a wallet holds several, which one answers
  a request for the document type?

  Source:
  \href{https://standards.atlassian.net/wiki/spaces/OneID/pages/2947416082/Technical+Implementation+Guide}{One
  ID guide} §4.1.2.
\end{definition}

\begin{definition}[Who issues a boarding pass (C-43)]
  \label{def:c43}
  \uses{def:boarding-pass, def:c01, def:c15}
  \notready
  The guide says a boarding pass is signed when it is issued and does not say
  who issues it or how a relying party trusts that issuer. Its trust wording
  is about issuing authority certificate lists.

  The issuer check of a boarding pass needs a list of boarding pass issuers,
  which may differ from the list of photoid issuers.

  Open point: Does each airline sign its own boarding passes, and how does an
  airport trust the airlines that fly there?

  Source:
  \href{https://standards.atlassian.net/wiki/spaces/OneID/pages/2947416082/Technical+Implementation+Guide}{One
  ID guide} §4.1.2.
\end{definition}

\begin{definition}[Which face image a touchpoint matches (C-97)]
  \label{def:c97}
  \uses{def:photoid, def:photoid-advisories}
  \notready
  A photoid may carry two face images: the portrait, which holds the
  passport's image, and the image captured at enrolment, which the guide says
  may be newer and of higher quality. The guide does not say which one a
  relying party matches a traveler's face against.

  A request names the data elements it wants, so the choice shows in the
  request. No rule says which is right.

  Open point: When a photoid carries both, which image does a relying party
  request and match: the portrait, the image captured at enrolment, or
  either?

  Source:
  \href{https://standards.atlassian.net/wiki/spaces/OneID/pages/2947416082/Technical+Implementation+Guide}{One
  ID guide} §4.1.1.
\end{definition}

\begin{definition}[A start of validity (C-98)]
  \label{def:c98}
  \uses{def:expired}
  \notready
  The guide names the end of a credential's validity and says nothing of a
  start.

  The model has no start of validity, so a credential is valid from any time
  until its validity ends.

  Open point: Does a credential have a start of validity that a relying party
  checks?

  Lookup: ISO/IEC 18013-5:2021 §9.3.1 and the validity information set when
  an mdoc is signed.

  Source:
  \href{https://standards.atlassian.net/wiki/spaces/OneID/pages/2947416082/Technical+Implementation+Guide}{One
  ID guide} §4.1.2.
\end{definition}

\begin{definition}[What a boarding pass string contains (C-102)]
  \label{def:c102}
  \uses{def:boarding-pass}
  \notready
  The guide refers the format of the boarding pass string, and which of its
  elements are mandatory, to IATA Resolution 792. It reproduces none of it.
  IATA's public guide to the resolution was read. It says one string holds
  one flight leg or several, that its items are mandatory, conditional or for
  an airline's own use, and that the passenger name is a mandatory item (BCBP
  guide §2.2.1). It leaves the list of items, their sizes and their values to
  the resolution's Attachments A and C (BCBP guide §2.2.1, §2.2.2).

  The model takes what the resolution decides as a parameter. With the
  resolution read, the passenger name, the legs and the mandatory elements of
  a string become definitions.

  Open point: Which items of a boarding pass string are mandatory, how long is
  the passenger name item and how is a longer name cut, and how are the legs
  of a journey and their dates given?

  Lookup: IATA Resolution 792, Attachments A and C.

  External: IATA Resolution 792 (not read); IATA BCBP Implementation Guide,
  seventh edition, §2.2.1, §2.2.2 (read).

  Source:
  \href{https://standards.atlassian.net/wiki/spaces/OneID/pages/2947416082/Technical+Implementation+Guide}{One
  ID guide} §4.1.2.
\end{definition}

\begin{definition}[What Tables C.1 to C.3 list (C-103)]
  \label{def:c103}
  \uses{def:photoid}
  \notready
  The guide refers the identifier, description, presence and encoding of
  every photoid data element to three tables of ISO/IEC TS 23220-4. It
  reproduces none of them.

  The model takes the tables as a parameter. With the tables read, they
  become one definition, and each guarantee proven for every table holds for
  it.

  Open point: Which data elements do the three tables list, which do they
  mark mandatory, and in which encoding?

  Lookup: ISO/IEC TS 23220-4:2026 Annex C, Tables C.1 to C.3.

  External: ISO/IEC TS 23220-4:2026 Annex C, Tables C.1 to C.3.

  Source:
  \href{https://standards.atlassian.net/wiki/spaces/OneID/pages/2947416082/Technical+Implementation+Guide}{One
  ID guide} §4.1.1.
\end{definition}

\begin{definition}[A boarding pass that covers an adult and an infant (C-104)]
  \label{def:c104}
  \uses{def:boarding-pass, def:traveler-match}
  \notready
  IATA's public guide to the bar code says some airlines issue one boarding
  pass for an adult travelling with an infant (BCBP guide §2.3.5). The One ID
  guide does not say whether a boarding pass credential may cover two
  travelers.

  The model gives a credential one subject. If one boarding pass covers two
  travelers, 'the same traveler' is not defined for it.

  Open point: May a boarding pass credential cover an adult and an infant? If
  it may, which of the two does it refer to when it is matched with a photoid,
  and what does the infant present?

  Blocks: G-03 (a boarding pass with two travelers has no one traveler to
  match).

  External: IATA BCBP Implementation Guide, seventh edition, §2.3.5 (read).

  Source:
  \href{https://standards.atlassian.net/wiki/spaces/OneID/pages/2947416082/Technical+Implementation+Guide}{One
  ID guide} §4.1.2.
\end{definition}

\begin{definition}[A relying party's check of a boarding pass's validity against the journey (C-109)]
  \label{def:c109}
  \uses{def:bp-validity}
  \notready
  The rule that a boarding pass's validity ends on or soon after the
  departure of the last leg is a duty of the issuer (C-38). No rule tells a
  relying party to check that the validity fits the journey. The only time
  check the guide gives a relying party is that a credential is within its
  validity (C-44).

  OID-104 is stated for an issuer. The check a relying party makes does not
  read which departure is meant, so for one credential its answer is the same
  under both readings. In the model the departure is an input of the issuer's
  rule and is not in the credential, so a relying party's check has nothing
  to compare the validity with.

  Case to check: A flight scheduled for 10:00 departs at 13:00, and the
  airline signs the boarding pass again so that its validity ends at 14:00.
  With ``soon after'' read as an hour, the validity breaks no requirement if
  the actual departure is meant, and ends long after the departure if the
  scheduled one is. At 12:30 the boarding pass is not expired under either
  reading, so a relying party that checks expiry accepts it in both.

  Open point: Must a relying party refuse a boarding pass whose validity does
  not fit the journey, or is expiry the only check?

  Source:
  \href{https://standards.atlassian.net/wiki/spaces/OneID/pages/2947416082/Technical+Implementation+Guide}{One
  ID guide} §4.1.2, consideration 1.
\end{definition}

% Candidate rules: functional requirements and assumptions (§3).  Written by hand; see docs/one-id-survey.md.

\section{Functional requirements and assumptions}

\begin{definition}[Requesting the minimum (C-11)]
  \label{def:c11}
  \uses{def:element}
  \notready
  A relying party must request the minimum disclosable attributes and verify
  them.

  A request is a set of data elements, and 'minimum' needs a purpose to be
  measured against. With one, the guarantee is that a request holds nothing
  its purpose does not need.

  Open point: What is the minimum for each use case: booking, advance
  passenger information, an e-gate, a boarding gate? §3.1 says the relying
  party must request the minimum, while §4.2.1.2 says it should request only
  what it needs, and §4.2.1.3 says the claims of a query must be the minimum.
  Which word holds?

  Source:
  \href{https://standards.atlassian.net/wiki/spaces/OneID/pages/2947416082/Technical+Implementation+Guide}{One
  ID guide} §3.1, item 3.
\end{definition}

\begin{definition}[Sharing in full or selectively (C-12)]
  \label{def:c12}
  \uses{def:c11, def:c07}
  \notready
  The passenger must share credentials securely, either in full or as the
  minimum selectively disclosable attributes.

  Both alternatives satisfy the sentence, so it does not limit what a wallet
  sends.

  Open point: May a wallet send a whole credential when the relying party
  asked for part of it?

  Source:
  \href{https://standards.atlassian.net/wiki/spaces/OneID/pages/2947416082/Technical+Implementation+Guide}{One
  ID guide} §3.1, item 4.
\end{definition}

\begin{definition}[Authenticity and integrity (C-13)]
  \label{def:c13}
  \uses{def:c01}
  What a relying party receives must be checked to be authentic and
  unaltered. §4.2 gives the checks for each profile.

  A response is accepted only when the issuer and device checks of its
  profile pass.

  Source:
  \href{https://standards.atlassian.net/wiki/spaces/OneID/pages/2947416082/Technical+Implementation+Guide}{One
  ID guide} §3.1, item 5.
\end{definition}

\begin{definition}[Security considerations (C-14)]
  \label{def:c14}
  \notready
  Security considerations must be followed for the request and the response
  of a presentation.

  Not a rule on its own. The model takes the security considerations to be
  the checks §4.2 lists.

  Open point: Which security considerations are meant: only those written in
  §4.2, or also those of the referenced specifications?

  Source:
  \href{https://standards.atlassian.net/wiki/spaces/OneID/pages/2947416082/Technical+Implementation+Guide}{One
  ID guide} §3.1, item 6.
\end{definition}

\begin{definition}[A list of trusted issuers (C-15)]
  \label{def:c15}
  \uses{def:c01}
  \notready
  The trust model is out of scope, but a relying party must have a mechanism
  to obtain the list of trusted issuers and intermediate certificate
  authorities.

  The list is an input of every issuer check.

  Open point: What is the mechanism, and what does a relying party do when
  the list cannot be obtained or is stale?

  Source:
  \href{https://standards.atlassian.net/wiki/spaces/OneID/pages/2947416082/Technical+Implementation+Guide}{One
  ID guide} §3.2, item 1.
\end{definition}

\begin{definition}[A list of trusted verifiers (C-16)]
  \label{def:c16}
  \uses{def:c01, def:c08}
  \notready
  Wallets must have access to a list of trusted verifiers or readers for
  authentication.

  If the list is required, a wallet refuses a request from a verifier that is
  not on it. The profiles do not say that.

  Open point: Is verifier authentication required? §3.2 says wallets must
  have the list. §4.2.1.2 says a wallet should be able to check verifier
  certificates, calls reader authentication recommended, and notes that not
  every wallet can process it. §4.2.2 says reader authentication should be
  present in a request. May a wallet answer a request it cannot authenticate?

  Source:
  \href{https://standards.atlassian.net/wiki/spaces/OneID/pages/2947416082/Technical+Implementation+Guide}{One
  ID guide} §3.2, item 2.
\end{definition}

\begin{definition}[Secure issuance (C-17)]
  \label{def:c17}
  \uses{def:c03}
  \notready
  Issuance of credentials to the wallet must be secure.

  An assumption with no rule in the guide behind it.

  Open point: What does 'securely issued' require? See C-03.

  Source:
  \href{https://standards.atlassian.net/wiki/spaces/OneID/pages/2947416082/Technical+Implementation+Guide}{One
  ID guide} §3.2, item 3.
\end{definition}

\begin{definition}[All credential types in the wallet (C-18)]
  \label{def:c18}
  \uses{def:photoid, def:boarding-pass, def:c42}
  \notready
  All credential types required for the passenger's journey must be present
  in the holder's wallet.

  A request for two document types can only be answered by a wallet that
  holds both.

  Open point: Which credential types does a journey require, and must they be
  in one wallet? A passenger with a photoid in one wallet and a boarding pass
  in an airline's app holds both and cannot present them together.

  Source:
  \href{https://standards.atlassian.net/wiki/spaces/OneID/pages/2947416082/Technical+Implementation+Guide}{One
  ID guide} §3.2, item 4.
\end{definition}

\begin{definition}[The holder is a passenger (C-19)]
  \label{def:c19}
  \uses{def:c03}
  \notready
  Whoever holds the wallet must be a passenger.

  Nothing in a credential says who holds the wallet. The model keeps the
  subject of each credential as a fact the issuer knows.

  Open point: Must every credential in a wallet be about its holder? A
  parent's wallet may hold the boarding passes of the children travelling
  with them. If that is allowed, holding a boarding pass does not make the
  holder its passenger.

  Source:
  \href{https://standards.atlassian.net/wiki/spaces/OneID/pages/2947416082/Technical+Implementation+Guide}{One
  ID guide} §3.2, item 5.
\end{definition}

\begin{definition}[Discovering contactless travel (C-20)]
  \label{def:c20}
  An airline may be able to find out which airports offer contactless travel.
  Annex B gives informative examples.

  Optional, and outside the credential and presentation rules.

  Source:
  \href{https://standards.atlassian.net/wiki/spaces/OneID/pages/2947416082/Technical+Implementation+Guide}{One
  ID guide} §3.2, item 6.
\end{definition}

\begin{definition}[Directing the passenger to the relying party (C-21)]
  \label{def:c21}
  \uses{def:c20}
  \notready
  Airlines must direct or prompt the passenger to the relying party's web
  page during the journey, if it is available.

  Outside the credential and presentation rules.

  Open point: What makes the page 'available', and what must an airline do
  when several relying parties serve one journey?

  Source:
  \href{https://standards.atlassian.net/wiki/spaces/OneID/pages/2947416082/Technical+Implementation+Guide}{One
  ID guide} §3.2, item 7.
\end{definition}

\begin{definition}[Authentication and consent before every presentation (C-22)]
  \label{def:c22}
  \uses{def:c07}
  \notready
  Every presentation of a credential must happen only after holder
  authentication and consent.

  A response exists only if the holder authenticated and consented to exactly
  the data elements in it.

  Open point: Is holder authentication required before every presentation?
  §3.2 says must. §4.2.1.3 says the response should be released after
  authentication, and §4.2.2 requires consent and does not mention
  authentication.

  Source:
  \href{https://standards.atlassian.net/wiki/spaces/OneID/pages/2947416082/Technical+Implementation+Guide}{One
  ID guide} §3.2, item 8.
\end{definition}

\begin{definition}[Terms and conditions (C-23)]
  \label{def:c23}
  \notready
  Terms and conditions must be explicitly mentioned to the passenger during
  the presentation process.

  Not modelled.

  Open point: Who shows the terms and conditions: the wallet, the relying
  party's page or the airline? Annex B mentions both the airline's and the
  relying party's.

  Source:
  \href{https://standards.atlassian.net/wiki/spaces/OneID/pages/2947416082/Technical+Implementation+Guide}{One
  ID guide} §3.2, item 9.
\end{definition}

\begin{definition}[Keys in a secure area (C-24)]
  \label{def:c24}
  \uses{def:c03}
  \notready
  Keys used for cryptography must be kept in a secure area.

  The model takes a device key to be usable only on the device it was issued
  to.

  Open point: Whose keys, and what counts as a secure area: hardware-backed
  storage on the phone, or any protected store? Does it apply to a verifier's
  request-signing and decryption keys?

  Source:
  \href{https://standards.atlassian.net/wiki/spaces/OneID/pages/2947416082/Technical+Implementation+Guide}{One
  ID guide} §3.2, item 10.
\end{definition}

% Candidate rules: remote presentation (§4.2.1).  Written by hand; see docs/one-id-survey.md.

\section{Remote presentation}

\begin{definition}[The order of a verifier's steps over the DC API (C-45)]
  \label{def:c45}
  A verifier must do four things in order: detect the API, check which
  protocols the browser supports, request the credentials by calling the
  browser's credential request, and validate the response.

  A response is validated only after a request the same verifier made.

  External: W3C Digital Credentials API §2.1 to §2.3.

  Source:
  \href{https://standards.atlassian.net/wiki/spaces/OneID/pages/2947416082/Technical+Implementation+Guide}{One
  ID guide} §4.2.1.1.
\end{definition}

\begin{definition}[Supporting both protocols (C-46)]
  \label{def:c46}
  \uses{def:c45}
  A verifier should support both protocols, \texttt{org-iso-mdoc} and
  \texttt{openid4vp-v1-signed}, and it is recommended to send both in one
  request when the browser supports both.

  A request may carry two protocol requests. The wallet answers one of them.

  External: W3C Digital Credentials API §2.2, §5.

  Source:
  \href{https://standards.atlassian.net/wiki/spaces/OneID/pages/2947416082/Technical+Implementation+Guide}{One
  ID guide} §4.2.1.1.
\end{definition}

\begin{definition}[Registering the wallet with the operating system (C-48)]
  \label{def:c48}
  \notready
  A wallet provider must register its app with the operating system's
  credential manager. The prerequisite is marked informative and uses MUST.

  Outside the model.

  Open point: Is registration a requirement or informative?

  Source:
  \href{https://standards.atlassian.net/wiki/spaces/OneID/pages/2947416082/Technical+Implementation+Guide}{One
  ID guide} §4.2.1.2, prerequisite 1 and step 2.
\end{definition}

\begin{definition}[The reader's list of trusted issuers (C-49)]
  \label{def:c49}
  \uses{def:c15}
  \notready
  A reader must have access to lists of trusted issuer certificates. Without
  them validation fails.

  The issuer check takes the list as an input.

  Open point: The same requirement as C-15: which list, and how is it
  obtained?

  Source:
  \href{https://standards.atlassian.net/wiki/spaces/OneID/pages/2947416082/Technical+Implementation+Guide}{One
  ID guide} §4.2.1.2, prerequisite 2.
\end{definition}

\begin{definition}[Cipher suites (C-50)]
  \label{def:c50}
  Reader and wallet must conform to the cipher suites of ISO/IEC 18013-5.

  Outside the model: cryptography is taken as ideal.

  External: ISO/IEC 18013-5:2021 §9.1.5.2.

  Source:
  \href{https://standards.atlassian.net/wiki/spaces/OneID/pages/2947416082/Technical+Implementation+Guide}{One
  ID guide} §4.2.1.2, prerequisite 3.
\end{definition}

\begin{definition}[The wallet's check of a verifier's certificate (C-51)]
  \label{def:c51}
  \uses{def:c16}
  \notready
  The heading says a wallet must have a trusted list of verifiers. The text
  says it should have a way to check the authority that issued a verifier's
  certificate against a trusted list.

  Whether a wallet refuses an unknown verifier depends on C-16.

  Open point: Heading and text differ: must or should?

  Source:
  \href{https://standards.atlassian.net/wiki/spaces/OneID/pages/2947416082/Technical+Implementation+Guide}{One
  ID guide} §4.2.1.2, prerequisite 4.
\end{definition}

\begin{definition}[Several document types in one request (C-53)]
  \label{def:c53}
  \uses{def:mdoc-request, def:photoid, def:boarding-pass}
  One request may name several document types, such as the photoid and the
  boarding pass.

  This is what 'presented together' means in the model: one response to one
  request that names both.

  Source:
  \href{https://standards.atlassian.net/wiki/spaces/OneID/pages/2947416082/Technical+Implementation+Guide}{One
  ID guide} §4.2.1.2, step 1.
\end{definition}

\begin{definition}[Requesting only what is needed (C-54)]
  \label{def:c54}
  \uses{def:c11, def:mdoc-request}
  \notready
  A relying party should request only the namespaces and data elements it
  needs.

  Same content as C-11 at a weaker level.

  Open point: §3.1 says must, this step says should. See C-11.

  Source:
  \href{https://standards.atlassian.net/wiki/spaces/OneID/pages/2947416082/Technical+Implementation+Guide}{One
  ID guide} §4.2.1.2, step 1.
\end{definition}

\begin{definition}[Reader authentication in the mdoc request (C-55)]
  \label{def:c55}
  \uses{def:c16, def:mdoc-request}
  \notready
  Reader authentication is recommended so the wallet can tell the request
  comes from a valid verifier. The guide notes not every wallet can process
  it, since it is optional in the specification, and says a verifier should
  take that into account.

  If a request need not be authenticated, a wallet may answer any verifier,
  and consent is the only gate.

  Open point: See C-16: is a request without reader authentication answered?

  External: ISO/IEC TS 23220-4:2026 Annex A §A.1.4.

  Source:
  \href{https://standards.atlassian.net/wiki/spaces/OneID/pages/2947416082/Technical+Implementation+Guide}{One
  ID guide} §4.2.1.2, step 1.
\end{definition}

\begin{definition}[The wallet's check of the request's origin (C-57)]
  \label{def:c57}
  \uses{def:c16, def:c55}
  \notready
  A wallet should check that a request comes from a trusted verifier, by
  validating reader authentication and the session transcript of the request.

  With 'should', a wallet that skips the check still conforms.

  Open point: See C-16.

  External: ISO/IEC 18013-5:2021 §9.1.4, §9.1.5.

  Source:
  \href{https://standards.atlassian.net/wiki/spaces/OneID/pages/2947416082/Technical+Implementation+Guide}{One
  ID guide} §4.2.1.2, step 3.
\end{definition}

\begin{definition}[The wallet's check of the request's content (C-58)]
  \label{def:c58}
  \uses{def:mdoc-request}
  A wallet should validate the request's syntax, document type, namespaces
  and data element identifiers, and may apply its own disclosure policies to
  accept or refuse the request.

  A wallet answers only for data elements it holds.

  Source:
  \href{https://standards.atlassian.net/wiki/spaces/OneID/pages/2947416082/Technical+Implementation+Guide}{One
  ID guide} §4.2.1.2, step 4.
\end{definition}

\begin{definition}[The mdoc response (C-59)]
  \label{def:c59}
  \uses{def:mdoc-request}
  The response must follow Annex C: it names the protocol \texttt{dcapi}, is
  encrypted to the verifier's key, and uses the session transcript of Annex
  C.

  A response is readable only by the verifier whose key the request carried,
  and its transcript ties it to that request.

  External: ISO/IEC TS 18013-7:2025 Annex C §C.3, §C.4, §C.5; IETF RFC 9180.

  Source:
  \href{https://standards.atlassian.net/wiki/spaces/OneID/pages/2947416082/Technical+Implementation+Guide}{One
  ID guide} §4.2.1.2, step 4.
\end{definition}

\begin{definition}[Decrypting the mdoc response (C-61)]
  \label{def:c61}
  \uses{def:c59}
  The verifier must decrypt the response with the mechanism of Annex C.

  Only the verifier that made the request can read the response.

  External: ISO/IEC TS 18013-7:2025 Annex C §C.4; IETF RFC 9180.

  Source:
  \href{https://standards.atlassian.net/wiki/spaces/OneID/pages/2947416082/Technical+Implementation+Guide}{One
  ID guide} §4.2.1.2, step 6.
\end{definition}

\begin{definition}[The passport security object check (C-64)]
  \label{def:c64}
  \uses{thm:portrait-is-dg2, def:photoid-advisories}
  \notready
  When a photoid carries data groups and the passport's security object, a
  verifier may check the security object as ICAO Doc 9303 Part 11 describes.

  The guide says the data groups namespace binds the credential to the
  passport chip's content. The check covers the data groups; whether it
  covers the portrait a touchpoint matches depends on the portrait being the
  same as data group 2.

  Open point: When is the check expected, and does it cover the portrait a
  touchpoint matches against?

  External: ICAO Doc 9303 Part 11 §5.

  Source:
  \href{https://standards.atlassian.net/wiki/spaces/OneID/pages/2947416082/Technical+Implementation+Guide}{One
  ID guide} §4.2.1.2 step 6, §4.2.1.3 step 5.
\end{definition}

\begin{definition}[What the mdoc profile does not check (C-65)]
  \label{def:c65}
  \uses{def:mdoc-request}
  \notready
  The verifier's list of checks in the mdoc profile has no step that compares
  the response with the request: that it holds only the document types and
  data elements requested, and that it answers this request's nonce. The
  OpenID4VP profile states a nonce check and a check that the token matches
  the query, and the proximity profile says issuer data validation covers the
  data elements requested.

  If the comparison is only implied by the ISO clauses, the model must assume
  it. If it is not, a response with more data than requested, or made for
  another request, is accepted.

  Case to check: A relying party requests the family name and receives a
  valid response that also holds the date of birth.

  Open point: Must a verifier in the mdoc profile check that the response
  matches its request and its nonce, as the OpenID4VP profile requires?

  Lookup: ISO/IEC TS 18013-7:2025 Annex C §C.3 to §C.5; ISO/IEC 18013-5:2021
  §9.1.3.

  Source:
  \href{https://standards.atlassian.net/wiki/spaces/OneID/pages/2947416082/Technical+Implementation+Guide}{One
  ID guide} §4.2.1.2, step 6.
\end{definition}

\begin{definition}[Reporting the result (C-66)]
  \label{def:c66}
  \uses{def:acceptance}
  \notready
  After processing a response, 'successful or error handled', the verifier
  must set the error structure that the ISO device response defines, and
  should return that standard's error codes.

  Not modelled.

  Open point: The structure cited is part of the wallet's response. What does
  a verifier set, and to whom does it return codes?

  External: ISO/IEC 18013-5:2021 §8.3.2.1.2.2, §8.3.2.1.2.3.

  Source:
  \href{https://standards.atlassian.net/wiki/spaces/OneID/pages/2947416082/Technical+Implementation+Guide}{One
  ID guide} §4.2.1.2 step 7, §4.2.2 step 6.
\end{definition}

\begin{definition}[Passing validated data on (C-67)]
  \label{def:c67}
  \uses{def:acceptance}
  After a successful validation the verifier should send the data to the
  identity management system through IATA's CUWS APIs.

  Data leaves the verifier only after acceptance.

  External: IATA CUWS standard APIs.

  Source:
  \href{https://standards.atlassian.net/wiki/spaces/OneID/pages/2947416082/Technical+Implementation+Guide}{One
  ID guide} §4.2.1.2 step 7, §4.2.1.3 step 6, §4.2.2 step 6.
\end{definition}

\begin{definition}[Complying with the high assurance profile (C-68)]
  \label{def:c68}
  \notready
  Verifier and wallet must comply with the OpenID4VC high assurance profile
  1.0 and follow Appendix A of OpenID4VP 1.0, and the profile's security and
  privacy considerations take precedence over the appendix's. Step 2 then
  says the request should comply with the profile's sections.

  Taken as must. For the DC API the profile's own considerations point to
  those of the appendix (HAIP §10, §11), so the precedence the guide states
  changes nothing.

  Open point: Must or should: the opening paragraph and step 2 differ.

  External: OpenID4VC HAIP 1.0 §5, §5.2, §5.3.1, §10, §11; OpenID4VP 1.0
  Appendix A (read).

  Source:
  \href{https://standards.atlassian.net/wiki/spaces/OneID/pages/2947416082/Technical+Implementation+Guide}{One
  ID guide} §4.2.1.3.
\end{definition}

\begin{definition}[The credential query (C-69)]
  \label{def:c69}
  \uses{def:element, def:c11}
  \notready
  The verifier must write a DCQL query: the credentials, each with the format
  \texttt{mso\_mdoc} and a document type, and the claims wanted, each a path
  of a namespace then a data element name. The claims must be the minimum
  needed. A claim may say whether the relying party intends to retain it;
  that member is optional.

  A query names data elements exactly as the mdoc request does, so one notion
  of request serves both profiles. In the guide's own first example one
  namespace is written with a leading space, which names no namespace of a
  photoid.

  Open point: What is the minimum set of claims? The same question as C-11,
  here stated with must.

  External: OpenID4VP 1.0 §6, §B.2.4.

  Source:
  \href{https://standards.atlassian.net/wiki/spaces/OneID/pages/2947416082/Technical+Implementation+Guide}{One
  ID guide} §4.2.1.3, step 1.
\end{definition}

\begin{definition}[Trusted authorities in the query (C-71)]
  \label{def:c71}
  \uses{def:c69, def:c15}
  Verifier and wallet must support the trusted authorities query, with which
  a verifier names the authorities whose credentials it accepts.

  A wallet can pick a credential the verifier will trust.

  External: OpenID4VP 1.0 §6.1.1.1; OpenID4VC HAIP 1.0 §5.

  Source:
  \href{https://standards.atlassian.net/wiki/spaces/OneID/pages/2947416082/Technical+Implementation+Guide}{One
  ID guide} §4.2.1.3, step 2.
\end{definition}

\begin{definition}[The nonce and state of the authorization request (C-72)]
  \label{def:c72}
  \notready
  The nonce and state must be set as OpenID4VP says. OpenID4VP requires a
  nonce that is fresh and random for every request, and asks for 'sufficient
  entropy' without a number (OpenID4VP §5.2). It does not define
  \texttt{state} over the DC API, where a wallet ignores it (OpenID4VP §5.3,
  §A.2). The guide's example request has a seven-digit nonce and no state,
  and the mdoc profile requires 16 bytes of entropy.

  A nonce is made fresh for every request, so a response made for one request
  does not answer another, provided the nonce cannot be guessed.

  Open point: What is the minimum entropy of the nonce in this profile: the 16
  bytes the mdoc profile requires? And is the must for \texttt{state} meant,
  given that it has no effect over the DC API?

  External: OpenID4VP 1.0 §5.2, §5.3, §A.2 (read).

  Source:
  \href{https://standards.atlassian.net/wiki/spaces/OneID/pages/2947416082/Technical+Implementation+Guide}{One
  ID guide} §4.2.1.3, step 2.
\end{definition}

\begin{definition}[Consent in the OpenID4VP profile (C-75)]
  \label{def:c75}
  \uses{def:c22, def:openid4vp-wallet-checks, def:openid4vp-binding}
  \notready
  Informed consent must be taken and the wallet must show what data is shared
  with whom. After user authentication and authorization the response
  'should' be released.

  Taken with §3.2: no response before authentication and consent.

  Open point: §3.2 says a presentation must happen only after authentication
  and consent. Does 'should be released' weaken that, or only describe the
  release?

  Source:
  \href{https://standards.atlassian.net/wiki/spaces/OneID/pages/2947416082/Technical+Implementation+Guide}{One
  ID guide} §4.2.1.3, step 4.
\end{definition}

\begin{definition}[Issuer and device checks in the OpenID4VP profile (C-78)]
  \label{def:c78}
  \uses{def:c15, def:c03, def:expired, def:c101}
  \notready
  The verifier must authenticate the issuer against a trusted list, which it
  'should' have access to, authenticate the issuer data, and authenticate the
  mdoc, which the guide calls holder binding. This profile's text does not
  name a time validity check (C-101).

  An accepted response comes from a trusted issuer and the holder's device.
  That it is not expired depends on C-101.

  Open point: §3.2 says a relying party must have the list of trusted
  issuers, this step says should. Which holds?

  External: ISO/IEC 18013-5:2021 §9.1.2, §9.1.3; ISO/IEC TS 23220-4:2026
  Annex A §A.3.1.

  Source:
  \href{https://standards.atlassian.net/wiki/spaces/OneID/pages/2947416082/Technical+Implementation+Guide}{One
  ID guide} §4.2.1.3, step 5.
\end{definition}

\begin{definition}[Errors in the OpenID4VP profile (C-79)]
  \label{def:c79}
  \uses{def:acceptance}
  Errors in processing a response should follow the OpenID4VP error
  responses.

  Not modelled.

  External: OpenID4VP 1.0 §8.5.

  Source:
  \href{https://standards.atlassian.net/wiki/spaces/OneID/pages/2947416082/Technical+Implementation+Guide}{One
  ID guide} §4.2.1.3, step 6.
\end{definition}

\begin{definition}[When the browser lacks the API or the protocol (C-99)]
  \label{def:c99}
  \uses{def:c45}
  \notready
  The verifier must check that the browser supports the API and the
  protocols. The guide does not say what the verifier does when it does not.

  A remote presentation has no defined outcome on a browser without support.

  Open point: What does a relying party offer a passenger whose browser or
  wallet supports neither protocol?

  Source:
  \href{https://standards.atlassian.net/wiki/spaces/OneID/pages/2947416082/Technical+Implementation+Guide}{One
  ID guide} §4.2.1.1.
\end{definition}

\begin{definition}[The version of the Digital Credentials API (C-100)]
  \label{def:c100}
  \uses{def:c45}
  \notready
  The guide puts a dated draft of the W3C Digital Credentials API in scope
  and notes the standard is still a draft.

  The protocol identifiers and the API are taken from that draft. A later
  draft may change them.

  Open point: When the W3C draft changes, which version must verifiers and
  wallets support, and who decides?

  External: W3C Digital Credentials API, draft of 16 July 2026.

  Source:
  \href{https://standards.atlassian.net/wiki/spaces/OneID/pages/2947416082/Technical+Implementation+Guide}{One
  ID guide} §1.4, §4.2.1.1.
\end{definition}

\begin{definition}[Whether issuer data authentication includes the validity check (C-101)]
  \label{def:c101}
  \uses{def:expired}
  \notready
  The mdoc profile cites one clause of ISO/IEC 18013-5 for issuer data
  authentication and says it covers time validity. The OpenID4VP profile and
  the proximity profile cite a different clause with must, and do not say
  whether it covers time validity; proximity adds the first clause with
  should.

  If the clause cited with must includes the validity check, an expired
  credential is refused in all three profiles. If it does not, two profiles
  accept one.

  Open point: Does the clause the OpenID4VP and proximity profiles cite with
  must include the time validity check?

  Lookup: ISO/IEC 18013-5:2021 §9.1.2 and §9.3.1.

  External: ISO/IEC 18013-5:2021 §9.1.2, §9.3.1.

  Source:
  \href{https://standards.atlassian.net/wiki/spaces/OneID/pages/2947416082/Technical+Implementation+Guide}{One
  ID guide} §4.2.1.3 step 5, §4.2.2 step 5.
\end{definition}

\begin{definition}[Checking the data model on what is received (C-107)]
  \label{def:c107}
  \uses{def:acceptance, def:photoid, def:boarding-pass}
  \notready
  The guide's lists of a relying party's checks hold no step that checks a
  response against the credential data model of §4.1: the mandatory data
  elements, their encodings, the format of the boarding pass string.  With
  selective disclosure a relying party sees only the data elements it
  requested.

  A relying party cannot tell from a response whether the credential behind
  it holds every mandatory data element.  The data model requirements can be
  checked in full only by the issuer, and by a wallet that holds the whole
  credential.

  Case to check: A relying party requests the family name of a photoid.  The
  photoid behind the response has no portrait, which one reading of the guide
  requires.  The response is the same as one from a complete photoid.

  Open point: Does a relying party check anything of §4.1 on a response, such
  as the encoding of what it received or the format of the boarding pass
  string?  Or are the data model requirements the issuer's duty alone?

  Source:
  \href{https://standards.atlassian.net/wiki/spaces/OneID/pages/2947416082/Technical+Implementation+Guide}{One ID guide} §4.1.1, §4.1.2, §4.2.1.3.
\end{definition}

% The presentation rules in Lean.

\begin{definition}[What a wallet and a relying party exchange]
  \label{def:exchange}
  \lean{Rules.OneId.Request, Rules.OneId.Consent, Rules.OneId.Response, Rules.OneId.Presented, Rules.OneId.Credential, Rules.OneId.ElementRef, Rules.OneId.Profile, Rules.OneId.SessionTranscript, Rules.OneId.DeviceSignature, Rules.OneId.Key, Rules.OneId.Nonce, Rules.OneId.Origin, Rules.OneId.Issuer, Rules.OneId.Mdoc.refs, Rules.OneId.Mdoc.disclose, Rules.OneId.Response.refs}
  \uses{def:mdoc, def:element}
  All three profiles have the same parts: a request that names data elements by
  document type, namespace and identifier, the holder's consent, a response
  that discloses data elements of credentials, and the relying party's checks.
  A credential as issued has an issuer and the key of the device it was issued
  to.  A presented credential has the data elements disclosed, the issuer's
  signature and a device signature.

  No cryptography is modelled.  Keys, nonces and origins are values that can
  only be compared, and a device signature is the key that made it and the
  session transcript it was made over.

  Source: \href{https://standards.atlassian.net/wiki/spaces/OneID/pages/2947416082/Technical+Implementation+Guide}{One ID guide} §4.2.1.1, §4.2.1.2, §4.2.1.3, §4.2.2.
\end{definition}

\begin{definition}[The session transcript of each profile]
  \label{def:session-transcript}
  \lean{Rules.OneId.IsoTranscripts, Rules.OneId.expectedTranscript}
  \uses{def:exchange, def:openid4vp-binding, def:c59}
  A device signs its response over a session transcript, and each profile
  defines the transcript differently.  The OpenID4VP profile's is in a public
  specification.  The other two are in ISO clauses this project has not read,
  so the model takes them as given: for a request, some value, with no
  assumption about which parts of the request it depends on.

  External: ISO/IEC TS 18013-7:2025 Annex C §C.5; ISO/IEC 18013-5:2021
  §9.1.5.1.

  Source: \href{https://standards.atlassian.net/wiki/spaces/OneID/pages/2947416082/Technical+Implementation+Guide}{One ID guide} §4.2.1.2, §4.2.2.
\end{definition}

\begin{definition}[A wallet's answer to a request (C-60)]
  \label{def:wallet-respond}
  \lean{Rules.OneId.Wallet.respond, Rules.OneId.Credential.present, Rules.OneId.ConsentReading, Rules.OneId.Request.wanted}
  \uses{def:exchange, def:session-transcript, def:c22}
  A wallet answers a request with the data elements that the request names and
  the holder approved, and with nothing of a credential that has none of them.
  It signs each credential it presents with that credential's device key, over
  the session transcript of the request as the wallet received it.

  The guide requires informed consent before every response, in all three
  profiles.  What a wallet does when the holder approves only some of the data
  elements requested is a reading: it sends those, or it sends nothing of that
  credential.

  External: OpenID4VP §6.4; ISO/IEC 18013-5:2021 §8.3.2.1.2.2, §9.1.3.

  Source: \href{https://standards.atlassian.net/wiki/spaces/OneID/pages/2947416082/Technical+Implementation+Guide}{One ID guide} §3.1, §3.2, §4.2.1.2, §4.2.1.3, §4.2.2.
\end{definition}

\begin{theorem}[A wallet discloses only what was requested and approved]
  \label{thm:wallet-discloses}
  \lean{Rules.OneId.wallet_discloses_only_requested_and_approved, Rules.OneId.wallet_presents_nothing_without_consent, Rules.OneId.wallet_signs_with_device_key}
  \uses{def:wallet-respond}
  Every data element a wallet discloses was named by the request and approved
  by the holder, under both readings of partial approval.  A wallet presents
  nothing when the holder approved nothing, and signs each credential it
  presents with that credential's device key over the transcript of the request
  it received.
\end{theorem}
\begin{proof}
  \uses{def:wallet-respond}
  The wallet discloses the requested data elements filtered by the approved ones.
\end{proof}

\begin{theorem}[Partial approval withholds the credential (C-105)]
  \label{thm:partial-consent}
  \lean{Rules.OneId.partial_consent_withholds_credential}
  \uses{def:wallet-respond}
  A wallet presents nothing of a credential when the holder approved only some
  of the data elements requested from it.  Not proven.

  Counterexample: a check-in page requests the family name and the portrait.
  The traveler approves the family name and declines the portrait.  A wallet
  that sends what was approved sends the family name.

  Open point: the guide requires informed consent and that the wallet show what
  is shared with whom, and does not say what happens when a holder agrees to a
  part.  OpenID4VP says a wallet that cannot deliver every claim asked of a
  credential must not return it (OpenID4VP §6.4.1).  When a holder approves
  only some of the data elements requested, does the wallet send those, or
  nothing from that credential?  Is the answer the same in all three profiles?

  Source: \href{https://standards.atlassian.net/wiki/spaces/OneID/pages/2947416082/Technical+Implementation+Guide}{One ID guide} §3.2, §4.2.1.3.
\end{theorem}

\begin{definition}[A relying party's decision on a response (C-47)]
  \label{def:acceptance}
  \lean{Rules.OneId.refusals, Rules.OneId.Presented.refusals, Rules.OneId.AcceptanceReading, Rules.OneId.AcceptanceReading.checksValidity, Rules.OneId.AcceptanceReading.checksTranscript, Rules.OneId.Refusal}
  \uses{def:exchange, def:session-transcript, def:expired, def:c03, def:c13, def:c15, def:c45, def:c49, def:c61, def:c65, def:c78, def:c88, def:c101}
  A relying party accepts a response only if every credential in it comes from
  an issuer on its list of trusted issuers, has an issuer's signature that
  verifies, and has a device signature made with the device key the issuer
  signed.  All three profiles require these.

  The profiles differ on three more checks, and the differences are readings:
  whether the device signature must be over the transcript of this request
  (required in the OpenID4VP profile), whether the end of validity is checked
  (required in the mdoc profile), and whether a response with a data element
  that was not requested is refused (stated nowhere).  Who is on the list of
  trusted issuers is given to the rule.  Decryption and revocation are not
  modelled.

  External: ISO/IEC 18013-5:2021 §9.1.2, §9.1.3, §9.3.1; ISO/IEC TS
  18013-7:2025 Annex C; OpenID4VP §8.6.

  Source: \href{https://standards.atlassian.net/wiki/spaces/OneID/pages/2947416082/Technical+Implementation+Guide}{One ID guide} §3.1, §3.2, §4.2.1.1, §4.2.1.2, §4.2.1.3, §4.2.2.
\end{definition}

\begin{theorem}[An accepted response comes from a trusted issuer]
  \label{thm:accepted-trusted-issuer}
  \lean{Rules.OneId.accepted_from_trusted_issuer}
  \uses{def:acceptance}
  Every credential in an accepted response comes from an issuer on the relying
  party's list and has an issuer's signature that verifies, in every profile
  and under every reading.
\end{theorem}
\begin{proof}
  \uses{def:acceptance}
  An issuer off the list and a signature that does not verify are each a reason to refuse.
\end{proof}

\begin{theorem}[An accepted response is signed with the device key]
  \label{thm:accepted-device-key}
  \lean{Rules.OneId.accepted_signed_by_device_key}
  \uses{def:acceptance}
  Every credential in an accepted response has a device signature made with the
  device key its issuer signed, in every profile and under every reading.  This
  is the device half of the guarantee G-02.
\end{theorem}
\begin{proof}
  \uses{def:acceptance}
  A missing device signature and one made with another key are each a reason to refuse.
\end{proof}

\begin{theorem}[A wallet cannot be checked for consent]
  \label{thm:no-check-shows-consent}
  \lean{Rules.OneId.no_check_shows_consent}
  \uses{def:acceptance, def:wallet-respond}
  Whatever a relying party checks on a request and its response, it cannot tell
  whether the holder approved what was sent.  The consent half of the guarantee
  G-01 rests on the wallet, and on whoever vouches for wallets.
\end{theorem}
\begin{proof}
  \uses{def:wallet-respond}
  The same response could have been sent with nothing approved.
\end{proof}

\begin{theorem}[An honest response is accepted]
  \label{thm:wallet-response-accepted}
  \lean{Rules.OneId.wallet_response_accepted}
  \uses{def:acceptance, def:wallet-respond}
  A response a wallet makes for the relying party's own request, from
  credentials of trusted issuers that have not expired, is accepted in every
  profile and under every reading.
\end{theorem}
\begin{proof}
  \uses{def:acceptance, def:wallet-respond}
  Each reason to refuse fails for such a response.
\end{proof}

\begin{definition}[A photoid and a boarding pass presented together (C-106)]
  \label{def:pair}
  \lean{Rules.OneId.pairRefusals, Rules.OneId.pairMatch, Rules.OneId.PairReading, Rules.OneId.PairRefusal, Rules.OneId.Presented.isOf}
  \uses{def:acceptance, def:traveler-match, def:c03, def:c53}
  A relying party accepts a photoid and a boarding pass presented together only
  if each is accepted on its own and the two pass a check that they are one
  traveler's.  The guide states no such check.  The content of the two can be
  compared, as in the credential data model, and a presentation adds the device
  keys: a reading in which the two credentials must carry the same device key.

  External: IATA Resolution 792; ISO/IEC 18013-5:2021 §9.1.3.

  Source: \href{https://standards.atlassian.net/wiki/spaces/OneID/pages/2947416082/Technical+Implementation+Guide}{One ID guide} §4.1.2, §3.2, §4.2.1.2.
\end{definition}

\begin{theorem}[In an accepted pair each credential is accepted]
  \label{thm:pair-each-accepted}
  \lean{Rules.OneId.accepted_pair_each_accepted}
  \uses{def:pair}
  In an accepted pair each credential is accepted on its own, the first is a
  photoid and the second is a boarding pass.
\end{theorem}
\begin{proof}
  \uses{def:pair}
  Each reason to refuse one credential is a reason to refuse the pair.
\end{proof}

\begin{theorem}[No check of a presentation identifies the traveler]
  \label{thm:no-presentation-check}
  \lean{Rules.OneId.no_presentation_check_identifies_the_traveler}
  \uses{def:pair, thm:no-content-check}
  Whatever a relying party checks on a presented photoid and a presented
  boarding pass, device keys included, a pair it accepts can be presentations
  of credentials issued about two different travelers.  A check of device keys
  gives one traveler only together with a rule about issuance.
\end{theorem}
\begin{proof}
  \uses{def:pair}
  The same two presentations could come from credentials issued about different people.
\end{proof}

\begin{theorem}[One traveler's pair is accepted]
  \label{thm:same-traveler-pair-accepted}
  \lean{Rules.OneId.same_traveler_pair_accepted}
  \uses{def:pair, thm:same-traveler-matches}
  A photoid and a boarding pass about the same traveler, each accepted on its
  own, are accepted together.  Not proven.

  Counterexample: a traveler holds a photoid on a phone and a boarding pass on
  a tablet.  Under the reading that the two credentials must carry the same
  device key, the pair is refused.

  Open point: if the device key is what ties the two credentials together, a
  traveler must hold both in one wallet on one device, and both issuers must
  issue to the same key.  Does One ID require a photoid and a boarding pass to
  be issued to the same device key, and what does a traveler do who holds them
  on two devices?

  Source: \href{https://standards.atlassian.net/wiki/spaces/OneID/pages/2947416082/Technical+Implementation+Guide}{One ID guide} §4.1.2, §3.2.
\end{theorem}

\begin{definition}[The OpenID4VP request (C-70)]
  \label{def:openid4vp-request}
  \lean{Rules.OneId.OpenId4Vp.requestFindings, Rules.OneId.OpenId4Vp.AuthorizationRequest, Rules.OneId.OpenId4Vp.CredentialQuery, Rules.OneId.OpenId4Vp.Certificate, Rules.OneId.OpenId4Vp.RequestReading, Rules.OneId.OpenId4Vp.RequestFinding, Rules.OneId.OpenId4Vp.AuthorizationRequest.shared, Rules.OneId.OpenId4Vp.protocolValue, Rules.OneId.OpenId4Vp.responseTypeValue, Rules.OneId.OpenId4Vp.responseModeValue, Rules.OneId.OpenId4Vp.clientIdPrefixValue, Rules.OneId.OpenId4Vp.mdocFormat}
  \uses{def:exchange, def:c45, def:c68, def:c69, def:c72}
  An OpenID4VP request is accepted only if it has the protocol value, response
  type, response mode and client identifier prefix the guide names, is signed
  with the key of the certificate its client identifier names, names its
  expected origins, carries an encryption key and a nonce, and asks only for
  mdoc credentials.  Whether it must also carry a state value is a reading.

  The rule does not check that the key and the nonce are new for each request,
  the minimum set of claims (C-69), or the entropy of the nonce (C-72).

  External: OpenID4VP §5, §5.1, §5.2, §5.9.3, §6, §A.2, §A.3.2; OpenID4VC HAIP
  §5, §5.2; W3C Digital Credentials API §5.

  Source: \href{https://standards.atlassian.net/wiki/spaces/OneID/pages/2947416082/Technical+Implementation+Guide}{One ID guide} §4.2.1.3.
\end{definition}

\begin{theorem}[An accepted OpenID4VP request has its required parts]
  \label{thm:openid4vp-request-parts}
  \lean{Rules.OneId.OpenId4Vp.accepted_request_parts}
  \uses{def:openid4vp-request}
  An accepted request has the four values the guide names, is signed with the
  key of the certificate its client identifier names, names an expected origin,
  carries an encryption key and a nonce, and asks only for mdoc credentials.
\end{theorem}
\begin{proof}
  \uses{def:openid4vp-request}
  Each missing part is a finding.
\end{proof}

\begin{theorem}[An accepted OpenID4VP request carries a state value]
  \label{thm:openid4vp-state}
  \lean{Rules.OneId.OpenId4Vp.accepted_request_has_state}
  \uses{def:openid4vp-request, def:c72}
  An accepted OpenID4VP request carries a state value.  Not proven.

  Counterexample: a signed request with every part the guide's step 2 lists and
  no state value, as in the guide's own example, is accepted under the reading
  that follows OpenID4VP.

  Open point: the guide says the nonce and state must be set as OpenID4VP §5.2
  says.  That clause requires state only for a presentation without a holder
  binding proof, §5.3 exempts the Digital Credentials API, and Appendix A, which
  the guide says must be followed, does not list the parameter.  Is the must
  for state meant, and if so what does a relying party or a wallet do with it?

  Source: \href{https://standards.atlassian.net/wiki/spaces/OneID/pages/2947416082/Technical+Implementation+Guide}{One ID guide} §4.2.1.3.
\end{theorem}

\begin{definition}[The wallet's checks of an OpenID4VP request (C-73)]
  \label{def:openid4vp-wallet-checks}
  \lean{Rules.OneId.OpenId4Vp.walletFindings, Rules.OneId.OpenId4Vp.WalletFinding}
  \uses{def:openid4vp-request, def:c16}
  A wallet answers an OpenID4VP request only if the request is well formed,
  came from an origin it names, is signed with the key of the certificate its
  client identifier names, and that certificate is a trusted relying party's.

  The guide introduces its list of four checks with should and marks three of
  them with must.  OpenID4VP, which the guide makes mandatory, requires the
  origin check and the check of the signature and the certificate chain of a
  wallet, so the model treats the checks as required.  Which list a wallet
  checks a relying party against is given to the rule (C-16).

  External: OpenID4VP §A.2, §5.9.3.

  Source: \href{https://standards.atlassian.net/wiki/spaces/OneID/pages/2947416082/Technical+Implementation+Guide}{One ID guide} §4.2.1.3, §3.2.
\end{definition}

\begin{theorem}[A wallet answers only a trusted, signed request from an expected origin]
  \label{thm:openid4vp-answered}
  \lean{Rules.OneId.OpenId4Vp.answered_request_parts}
  \uses{def:openid4vp-wallet-checks}
  A request a wallet answers came from an origin it names, is signed with the
  key of the certificate its client identifier names, and that certificate is
  on the wallet's list of trusted relying parties.
\end{theorem}
\begin{proof}
  \uses{def:openid4vp-wallet-checks}
  Each failed check is a finding.
\end{proof}

\begin{definition}[What an OpenID4VP response is bound to (C-77)]
  \label{def:openid4vp-binding}
  \lean{Rules.OneId.OpenId4Vp.transcript, Rules.OneId.OpenId4Vp.deviceResponses}
  \uses{def:exchange}
  In the OpenID4VP profile the wallet signs its response over the origin of the
  request, the nonce and the relying party's encryption key, and the relying
  party rebuilds that transcript and validates the nonce.  Each mdoc is
  returned in a device response of its own.

  The specification hashes the three values.  The model keeps the values, so it
  does not model two inputs with one hash.

  External: OpenID4VP §B.2.5, §B.2.6.2; OpenID4VC HAIP §5.3.1.

  Source: \href{https://standards.atlassian.net/wiki/spaces/OneID/pages/2947416082/Technical+Implementation+Guide}{One ID guide} §4.2.1.3.
\end{definition}

\begin{theorem}[An OpenID4VP response cannot be used for another request]
  \label{thm:openid4vp-refuses-other-request}
  \lean{Rules.OneId.OpenId4Vp.refuses_other_request, Rules.OneId.OpenId4Vp.accepted_for_one_request, Rules.OneId.OpenId4Vp.one_mdoc_per_deviceResponse}
  \uses{def:openid4vp-binding, def:acceptance}
  In the OpenID4VP profile a response signed for another nonce, another origin
  or another encryption key is refused, under every reading.  A response
  accepted for two requests was made for one nonce, so a captured response
  cannot answer a later request.
\end{theorem}
\begin{proof}
  \uses{def:openid4vp-binding, def:acceptance}
  The transcript the relying party rebuilds differs when any of the three values differs.
\end{proof}

\begin{definition}[The request of the mdoc profile (C-52)]
  \label{def:mdoc-request}
  \lean{Rules.OneId.MdocProfile.requestFindings, Rules.OneId.MdocProfile.DeviceRequest, Rules.OneId.MdocProfile.RequestFinding, Rules.OneId.MdocProfile.DeviceRequest.shared, Rules.OneId.MdocProfile.protocolValue, Rules.OneId.MdocProfile.encryptionProtocolValue, Rules.OneId.MdocProfile.minimumNonceEntropy}
  \uses{def:exchange, def:c45, def:element}
  A request of the mdoc profile is accepted only if it has the protocol value
  \texttt{org-iso-mdoc}, its encryption information names \texttt{dcapi}, and it
  carries a nonce of at least 16 bytes of entropy and the recipient's public
  key.  The entropy of a nonce is taken as given.  Reader authentication is
  recommended and is not checked by this rule.

  External: ISO/IEC TS 18013-7:2025 Annex C §C.2; W3C Digital Credentials API
  §5.

  Source: \href{https://standards.atlassian.net/wiki/spaces/OneID/pages/2947416082/Technical+Implementation+Guide}{One ID guide} §4.2.1.2.
\end{definition}

\begin{theorem}[An accepted mdoc request has its required parts]
  \label{thm:mdoc-request-parts}
  \lean{Rules.OneId.MdocProfile.accepted_deviceRequest_parts}
  \uses{def:mdoc-request}
  An accepted request of the mdoc profile names both protocol values and
  carries a nonce of at least 16 bytes of entropy and a recipient key.
\end{theorem}
\begin{proof}
  \uses{def:mdoc-request}
  Each missing part is a finding.
\end{proof}

% Candidate rules: proximity presentation (§4.2.2).  Written by hand; see docs/one-id-survey.md.

\section{Proximity presentation}

\begin{definition}[When proximity presentation is used (C-80)]
  \label{def:c80}
  \uses{thm:photoid-has-portrait}
  Proximity presentation is required for one-to-one face matching when the
  face image was not shared in advance, and should be used whenever
  credentials are shared on demand or offline.

  The touchpoint needs a portrait from the photoid at that moment, so the
  profile depends on the portrait being present.

  Source:
  \href{https://standards.atlassian.net/wiki/spaces/OneID/pages/2947416082/Technical+Implementation+Guide}{One
  ID guide} §4.2.2.
\end{definition}

\begin{definition}[The proximity request (C-84)]
  \label{def:c84}
  \uses{def:proximity-session, def:element}
  The reader must send a device request as ISO/IEC 18013-5 defines it,
  encrypted with the session key. Reader authentication should be present,
  and when it is its structure must follow the same standard.

  A request lists a document type, namespaces and data elements, as in the
  remote profiles.

  External: ISO/IEC 18013-5:2021 §8.3.2.1.2.1, §9.1.4.

  Source:
  \href{https://standards.atlassian.net/wiki/spaces/OneID/pages/2947416082/Technical+Implementation+Guide}{One
  ID guide} §4.2.2, step 3.
\end{definition}

\begin{definition}[The wallet's checks in proximity (C-85)]
  \label{def:c85}
  \uses{def:c84, def:c16}
  \notready
  After decrypting, the wallet must check that the document type matches a
  credential it holds, that the requested data elements match what is stored,
  and must perform reader authentication and validate the session transcript.

  The wallet must authenticate the reader, while a request need only 'should'
  carry reader authentication (C-84).

  Open point: What does a wallet do with a proximity request that has no
  reader authentication: refuse it, or answer it?

  External: ISO/IEC 18013-5:2021 §9.1.4, §9.1.5.

  Source:
  \href{https://standards.atlassian.net/wiki/spaces/OneID/pages/2947416082/Technical+Implementation+Guide}{One
  ID guide} §4.2.2, step 4.
\end{definition}

\begin{definition}[Decrypting the proximity response (C-88)]
  \label{def:c88}
  \uses{def:wallet-respond}
  \notready
  The reader 'should' decrypt the response with the session keys and rebuild
  the session transcript.

  Taken as must: nothing can be validated otherwise.

  Open point: 'Should' for a step every later check needs. Is it must?

  External: ISO/IEC 18013-5:2021 §9.1.5, §9.1.5.1.

  Source:
  \href{https://standards.atlassian.net/wiki/spaces/OneID/pages/2947416082/Technical+Implementation+Guide}{One
  ID guide} §4.2.2, step 5.
\end{definition}

\begin{definition}[Ending the session (C-91)]
  \label{def:c91}
  \uses{def:proximity-session}
  The session shall be terminated when the steps are done or on a condition
  such as a timeout, and all session data made for it must be deleted.

  A session's keys are not used again.

  External: ISO/IEC 18013-5:2021 §9.1.1.4.

  Source:
  \href{https://standards.atlassian.net/wiki/spaces/OneID/pages/2947416082/Technical+Implementation+Guide}{One
  ID guide} §4.2.2, step 7.
\end{definition}

% The proximity rules in Lean.

\begin{definition}[Keys and retrieval in a proximity session (C-82)]
  \label{def:proximity-session}
  \lean{Rules.OneId.Proximity.sessionFindings, Rules.OneId.Proximity.Session, Rules.OneId.Proximity.Retrieval, Rules.OneId.Proximity.SessionFinding, Rules.OneId.Proximity.Session.shared}
  \uses{def:exchange}
  The device makes a key pair for each engagement, the reader makes one before
  each engagement, and data is retrieved from the device and never from a
  server.  The rule looks at the sessions of one deployment together, since a
  reused key shows only across sessions.  The end of a session and the
  deletion of its data are not modelled.

  External: ISO/IEC 18013-5:2021 §8.2.1.1, §9.1.1.4, §9.1.1.5.

  Source: \href{https://standards.atlassian.net/wiki/spaces/OneID/pages/2947416082/Technical+Implementation+Guide}{One ID guide} §4.2.2.
\end{definition}

\begin{theorem}[Accepted proximity sessions use fresh keys and no server]
  \label{thm:proximity-sessions}
  \lean{Rules.OneId.Proximity.accepted_sessions}
  \uses{def:proximity-session}
  In accepted proximity sessions no session retrieves data from a server, no
  reader's key is used twice and no device's engagement key is used twice.
\end{theorem}
\begin{proof}
  \uses{def:proximity-session}
  Server retrieval and a repeated key are each a finding.
\end{proof}

% Candidate rules: what the guide leaves out (§1.4 and §2).  Written by hand; see docs/one-id-survey.md.

\section{What the guide leaves out}

\begin{definition}[The trust model (C-01)]
  \label{def:c01}
  \notready
  The guide puts the trust model out of scope: how trust is anchored,
  delegated and verified between issuer, wallet and verifier. It still relies
  on lists of trusted issuers and trusted verifiers throughout §3.2 and §4.2.

  Every check that says 'trusted' takes the list as an input the guide does
  not define. A guarantee about authenticity holds only relative to that
  list.

  Open point: Who accredits an issuer of photoids, and who accredits an
  issuer of boarding passes? Is it one list or two, who publishes it, and how
  does a relying party or a wallet obtain it and know it is current?

  Blocks: G-01 (the wallet cannot say whom data is shared with unless
  verifiers are on a list it trusts), G-02 (a device key proves holder
  binding only if the issuer that signed it is trusted), G-03 (a boarding
  pass and a photoid come from different issuers, and matching them assumes
  both are trusted), G-04 (the validity a verifier checks is the one a
  trusted issuer signed).

  Source:
  \href{https://standards.atlassian.net/wiki/spaces/OneID/pages/2947416082/Technical+Implementation+Guide}{One
  ID guide} §1.4, §2.
\end{definition}

\begin{definition}[Revocation (C-02)]
  \label{def:c02}
  \notready
  The guide puts revocation out of scope. Within its rules a credential stops
  being accepted when the validity its issuer signed ends, or when its issuer
  is no longer trusted.

  No rule can refuse a credential that is still within its validity. A
  guarantee that an expired credential is refused says nothing about a
  cancelled boarding pass or a revoked passport.

  Open point: How does a boarding pass credential stop being accepted when
  the passenger is rebooked, offloaded or cancels, and a photoid when the
  passport is reported lost or the phone is stolen? Is short validity the
  intended answer, and if so how short?

  Blocks: G-04 (an expired credential is refused, a withdrawn one is not),
  G-03 (a boarding pass that was reissued to a rebooked passenger still
  matches the photoid), G-02 (a device key on a stolen phone stays bound).

  Source:
  \href{https://standards.atlassian.net/wiki/spaces/OneID/pages/2947416082/Technical+Implementation+Guide}{One
  ID guide} §1.4.
\end{definition}

\begin{definition}[Issuance to the wallet (C-03)]
  \label{def:c03}
  \notready
  The guide puts normative issuance out of scope and gives only an
  informative annex. It assumes credentials are securely issued (§3.2). The
  rules it does give that are fixed at issuance are the content of each
  credential (§4.1.1, §4.1.2) and the end of validity of a boarding pass.

  Facts fixed at issuance, such as which device key a credential is bound to
  and who its subject is, enter the model as inputs with no rule behind them.

  Open point: Which issuance requirements does a relying party depend on? At
  the least: that the device key is bound to the holder's device, that the
  issuer checked the holder against the passport, and that a boarding pass is
  issued to the passenger it names.

  Blocks: G-02 (binding to the device key starts at issuance), G-03 (the
  subject of each credential is set at issuance).

  Source:
  \href{https://standards.atlassian.net/wiki/spaces/OneID/pages/2947416082/Technical+Implementation+Guide}{One
  ID guide} §1.4.
\end{definition}

\begin{definition}[The data elements of a photoid (C-04)]
  \label{def:c04}
  \uses{def:c103}
  \notready
  The guide puts the data modelling and schema of credentials out of scope,
  except for the boarding pass. The data elements of a photoid are in tables
  of ISO/IEC TS 23220-4 that the guide refers to.

  The photoid rules take the three tables as a parameter. Each guarantee is
  proven for every content the tables could have, or is open.

  Open point: Which data elements does One ID require in a photoid beyond
  what the ISO tables require? The guide names only the portrait, and its
  examples request the names and an age attestation.

  External: ISO/IEC TS 23220-4:2026 Annex C, Tables C.1 to C.3.

  Source:
  \href{https://standards.atlassian.net/wiki/spaces/OneID/pages/2947416082/Technical+Implementation+Guide}{One
  ID guide} §1.4.
\end{definition}

\begin{definition}[User interface and user experience (C-05)]
  \label{def:c05}
  \notready
  The guide puts user interface and experience out of scope, and also
  requires wallets to show what data is shared with whom (§4.2.1.2, §4.2.1.3)
  and to display a boarding pass in a user-friendly way (§4.1.2).

  The display obligations cannot be stated precisely while the interface is
  out of scope. The model treats consent as a recorded decision about a set
  of data elements and a named verifier.

  Open point: Are the display obligations in §4.1.2 and §4.2 requirements a
  wallet is tested against, or guidance?

  Source:
  \href{https://standards.atlassian.net/wiki/spaces/OneID/pages/2947416082/Technical+Implementation+Guide}{One
  ID guide} §1.4.
\end{definition}

\begin{definition}[What a presentation may contain (C-07)]
  \label{def:c07}
  \uses{def:element}
  A presentation may include selectively disclosed claims or all claims, from
  one or more credentials.

  A presentation is a set of data elements drawn from one or more
  credentials. Needed by every guarantee about disclosure.

  Source:
  \href{https://standards.atlassian.net/wiki/spaces/OneID/pages/2947416082/Technical+Implementation+Guide}{One
  ID guide} §2.
\end{definition}

\begin{definition}[Who a trust list covers (C-08)]
  \label{def:c08}
  \uses{def:c01}
  \notready
  A trust list publishes accredited issuers and may be extended to verifiers
  and wallet providers.

  Whether a list of trusted verifiers exists at all is left open here, while
  §3.2 requires wallets to have one.

  Open point: Does the One ID trust list cover verifiers and wallet
  providers, or issuers only?

  Source:
  \href{https://standards.atlassian.net/wiki/spaces/OneID/pages/2947416082/Technical+Implementation+Guide}{One
  ID guide} §2.
\end{definition}

% Candidate rules: issuance and contactless travel (Annexes A and B, informative).  Written by hand; see docs/one-id-survey.md.

\section{Issuance and contactless travel}

\begin{definition}[The issuance protocol (C-92)]
  \label{def:c92}
  \uses{def:c03}
  Issuance with OpenID4VCI 1.0, profiled by the high assurance profile, is
  recommended. The annex is informative.

  Does not fill the gap of C-03: nothing here is a requirement.

  External: OpenID4VCI 1.0; OpenID4VC HAIP 1.0 §4.

  Source:
  \href{https://standards.atlassian.net/wiki/spaces/OneID/pages/2947416082/Technical+Implementation+Guide}{One
  ID guide} Annex A, §5.1.
\end{definition}

\begin{definition}[A boarding pass string sent by the airline (C-93)]
  \label{def:c93}
  \uses{thm:matched-pair-same-traveler, def:boarding-pass}
  \notready
  In the informative flow an airline may send the boarding pass string to the
  relying party's API itself, with the time consent was recorded, before the
  passenger presents anything.

  A boarding pass string that arrives this way is not signed by an issuer and
  not bound to a device. It is matched to the photoid by the session alone.

  Open point: When the airline sends the boarding pass string, what ties it
  to the photoid the passenger presents later, and whose consent does the
  recorded time stand for?

  Source:
  \href{https://standards.atlassian.net/wiki/spaces/OneID/pages/2947416082/Technical+Implementation+Guide}{One
  ID guide} Annex B, §6.2.1, step 2.
\end{definition}

\begin{definition}[What a contactless relying party requests (C-94)]
  \label{def:c94}
  \uses{def:c11, def:c53, thm:photoid-has-portrait, thm:matched-pair-same-traveler, def:c93}
  The relying party requests the boarding pass and selected photoid data
  elements including the portrait, or only the photoid data elements when the
  airline already sent the boarding pass string.

  The one place the guide says what a minimum request is for a use case: the
  portrait is always in it.

  Source:
  \href{https://standards.atlassian.net/wiki/spaces/OneID/pages/2947416082/Technical+Implementation+Guide}{One
  ID guide} Annex B, §6.2.1, step 5.
\end{definition}

\begin{definition}[The session nonce of the relying party's API (C-95)]
  \label{def:c95}
  \uses{def:c56, def:c72}
  \notready
  The relying party's API returns a session identifier, a nonce unique to the
  session, issue and expiry times, and the address of its page.

  The annex does not say whether this nonce is the one in the presentation
  request.

  Open point: Is the session nonce the presentation nonce? If not, what ties
  the presentation to the airline's session?

  Source:
  \href{https://standards.atlassian.net/wiki/spaces/OneID/pages/2947416082/Technical+Implementation+Guide}{One
  ID guide} Annex B, §6.2.1, step 3.
\end{definition}

\begin{definition}[Several relying parties in one journey (C-96)]
  \label{def:c96}
  \uses{def:c21, def:c22}
  When a journey has several relying parties, the passenger presents to each
  in turn.

  Each presentation is its own request with its own consent.

  Source:
  \href{https://standards.atlassian.net/wiki/spaces/OneID/pages/2947416082/Technical+Implementation+Guide}{One
  ID guide} Annex B, §6.2.1, step 7.
\end{definition}

% Candidate rules: candidate guarantees (the four guarantees this project was asked to check).  Written by hand; see docs/one-id-survey.md.

\section{Candidate guarantees}

\begin{theorem}[Only requested and consented data is received (G-01)]
  \label{thm:g01}
  \lean{Rules.OneId.accepted_holds_only_requested}
  \uses{def:acceptance, thm:wallet-discloses, thm:no-check-shows-consent, thm:partial-consent, def:c26, def:c07, def:c12, def:c22, def:c58, def:c65, def:c75, def:element, thm:names-unique, def:c31}
  A relying party receives only the data elements it requested and the holder
  consented to.  Not proven.

  A wallet that follows the rules discloses only what was requested and
  approved, which is proven.  Whether an accepted response holds only what was
  requested depends on the relying party refusing extra data, which no text
  that was read requires.  And no check a relying party makes can show that
  the holder approved what was sent, so that half rests on the wallet.

  Counterexample: a check-in page requests only the family name.  A wallet
  that does not follow the rules also sends the portrait, with valid
  signatures made for this request.  A relying party that does not refuse
  extra data accepts the response and holds a face image it never asked for.

  Open point: OpenID4VP forbids a wallet to send a claim the query did not
  select and tells a verifier not to rely on the wallet for it, without
  saying that a response with more is refused (OpenID4VP §6.4, §14.9).  The
  mdoc profile lists no step that compares the response with the request
  (C-65).  Must a relying party refuse a response that holds a data element
  it did not request, or may it accept the response and discard the extra?
  And who vouches that a wallet asks for consent (C-01, C-08)?

  Source:
  \href{https://standards.atlassian.net/wiki/spaces/OneID/pages/2947416082/Technical+Implementation+Guide}{One ID guide} §4.1.1, §3.1, §4.2.1.3.
\end{theorem}

\begin{theorem}[An accepted presentation is bound to the device and the request (G-02)]
  \label{thm:g02}
  \lean{Rules.OneId.accepted_answers_this_request}
  \uses{def:acceptance, thm:accepted-device-key, thm:openid4vp-refuses-other-request, def:c78, def:c65, def:c72, def:c59, def:c88, def:proximity-session, def:c03, def:c24, def:c02}
  An accepted presentation is bound to the holder's device key and to the
  verifier's request: its nonce in the remote profiles, its session in
  proximity.  Not proven.

  The device half is proven for all three profiles: every accepted credential
  is signed with the device key its issuer signed.  The request half is proven
  for the OpenID4VP profile, where a response made for one nonce, origin and
  key is refused for any other.  It is open for the mdoc profile and for
  proximity.

  Counterexample: a response made for a request with one nonce is captured and
  sent again in answer to a later request with another nonce.  The OpenID4VP
  profile refuses it.  The mdoc profile accepts it if the relying party does
  not compare the session transcript with its own request.

  Open point: the OpenID4VP profile requires the relying party to rebuild the
  session transcript and validate the nonce.  The mdoc profile's list of
  checks has no such step (C-65), and the proximity profile lists the
  rebuilding with no requirement word (C-88).  Must a relying party in the
  mdoc profile and a reader in proximity check that a response was made for
  their own request?  The guarantee also rests on issuance for what a device
  key is bound to (C-03), and on the entropy of the nonce (C-72).

  Source:
  \href{https://standards.atlassian.net/wiki/spaces/OneID/pages/2947416082/Technical+Implementation+Guide}{One ID guide} §4.2.
\end{theorem}

\begin{theorem}[A photoid and a boarding pass presented together are one traveler's (G-03)]
  \label{thm:g03}
  \lean{Rules.OneId.accepted_pair_same_traveler}
  \uses{def:pair, thm:no-presentation-check, thm:matched-pair-same-traveler, thm:same-traveler-matches, def:c19, def:c42, def:c93, def:c94, def:c01, def:c02, def:c03, def:c104}
  A photoid and a boarding pass presented together refer to the same
  traveler, under a defined matching rule.  Not proven.

  It fails under every reading of the check: with nothing compared, with
  names compared, and with device keys compared.  Comparing device keys gives
  one traveler where issuance gives each device key to one traveler, which is
  proven under that condition and which the guide does not state.

  Counterexample: Alice Martin's photoid and Bruno Keller's boarding pass were
  issued to one shared phone.  Under the reading that the two credentials must
  carry the same device key, the pair is accepted: Alice's face goes with
  Bruno's boarding pass.

  Open point: a presentation gives a relying party one more thing to compare
  than the content of the two credentials: the device key each was issued to.
  That settles the question only if issuance makes a device key stand for one
  traveler, and the guide puts issuance out of scope (C-03).  What ties a
  photoid and a boarding pass presented together to one traveler: the same
  device key, a rule for issuers, a comparison of names, or something else,
  and what must a relying party check?  A boarding pass for an adult with an
  infant is a further case (C-104).

  Source:
  \href{https://standards.atlassian.net/wiki/spaces/OneID/pages/2947416082/Technical+Implementation+Guide}{One ID guide} §4.1.2, §3.2.
\end{theorem}

\begin{theorem}[An expired credential is never accepted (G-04)]
  \label{thm:g04}
  \lean{Rules.OneId.accepted_not_expired}
  \uses{def:acceptance, def:expired, def:c78, def:c101, thm:valid-at-validuntil, def:c02}
  A relying party never accepts a credential whose validity has ended.  Not
  proven.

  It is proven for the mdoc profile, where the guide says the clause it cites
  covers time validity.  It is open for the OpenID4VP profile and for
  proximity.  Within the guide's rules a credential from a trusted issuer
  ends only by expiry, since revocation is out of scope (C-02).

  Counterexample: a boarding pass whose validity ended at 10:00 is presented
  at 11:00 with valid signatures.  The mdoc profile refuses it.  The OpenID4VP
  profile and proximity accept it if the clause they cite for issuer data
  authentication has no validity check.

  Open point: the OpenID4VP and proximity profiles cite a different clause of
  ISO/IEC 18013-5 from the mdoc profile, with must, and do not say whether it
  covers time validity (C-101).  Is a relying party required to refuse an
  expired credential in all three profiles?

  Lookup: ISO/IEC 18013-5:2021 §9.1.2 and §9.3.1.

  Source:
  \href{https://standards.atlassian.net/wiki/spaces/OneID/pages/2947416082/Technical+Implementation+Guide}{One ID guide} §4.1.2, §4.2.
\end{theorem}


\input{ears.generated}
